\section{A quantitative reduction of the ramified global cutoff problem}
\label{sec:ACcutoffbridge}

This section continues the attempt to prove T-CUT from AC-T11.
The calculations below use the actual theta images and their
orthogonal projections.  They distinguish the scalar prolate
problem from the additional arithmetic norm.
The new estimates in AC-T28--AC-T30 now prove subpower
arithmetic inverse-Gram control and the actual tested
comparison with a positive norm-window contraction at
C99's full cutoff.  AC-T31 proves the full-multiplicity
zero-measure lower bound.  AC-T32--AC-T34 identify the
actual co-Poisson adjoint and prove a uniform raw nuclear
bound with compact boundary limits.  The remaining
identification is the matching upper limit in (TB.53),
including control after arithmetic polar normalization.
Fix a primitive nonprincipal Dirichlet character \(\chi\) of
conductor \(q\), with \(\chi(-1)=(-1)^\epsilon\).
For \(\Lambda>1\), let \(\psi_1,\ldots,\psi_N\) be the prolate
vectors of parity \(\epsilon\) chosen in AC-T11, supported in
\([-\Lambda,\Lambda]\), with the conductor-adjusted frequency
parameter \(\Lambda/q\).  The intended full family has indices
up to \(4\Lambda^2/q\).
Put
\[
 e_j(x)=E_\chi\psi_j(x)
   =\sqrt{x}\sum_{n\geq1}\chi(n)\psi_j(nx),\quad
 G_{ij}=\langle e_i,e_j\rangle,\quad
 S_\Lambda=1_{[\Lambda^{-1},\Lambda]}.
\]
Inner products are conjugate-linear in the first variable.
The Hilbert space is \(L^2(\mathbb R_{>0},dx/x)\).

\begin{proposition}[AC-T21: the exact arithmetic Gram weight]
\label{AC-T21}
Define
\[
 m_j(\tau)=\int_0^\infty
       \psi_j(x)x^{-1/2+i\tau}\,dx,\qquad
 m(\tau)=(m_1(\tau),\ldots,m_N(\tau)).
\]
Then
\[
 G=\frac1{2\pi}\int_{\mathbb R}
      |L(1/2+i\tau,\chi)|^2 m(\tau)^*m(\tau)\,d\tau,
 \tag{TB.1}
\]
and, for the orthogonal projection \(Q_\Lambda\) onto
\(\operatorname{span}\{e_j\}\),
\[
 \operatorname{Tr}(Q_\Lambda V(h))
   =\frac1{2\pi}\int_{\mathbb R}\widehat h_+(\tau)
       K_{\chi,\Lambda}(\tau)\,d\tau,\qquad
 K_{\chi,\Lambda}(\tau)=
 |L(1/2+i\tau,\chi)|^2
           m(\tau)G^{-1}m(\tau)^* .
 \tag{TB.2}
\]
These integrals are absolutely convergent for bounded
\(\widehat h_+\).

Conjugating the original time-frequency rectangle by dilation
\(\sqrt q\) replaces the scalar prolate parameter by
\(\Lambda/\sqrt q\).  This dilation multiplies every \(m_j(\tau)\)
by the same phase \(q^{i\tau/2}\), and so leaves (TB.1)--(TB.2)
unchanged apart from the replacement of the scalar prolate basis.
The weight \(|L(1/2+i\tau,\chi)|^2\) remains.
For nonprincipal \(\chi\), it is not a constant multiple of
\(|\zeta(1/2+i\tau)|^2\).
\end{proposition}
\begin{proof}
Let \(B_\chi=\sup_M|\sum_{n\leq M}\chi(n)|<\infty\).
Summation by parts bounds
\(\sum\chi(n)\psi_j(nx)\) uniformly in \(x>0\) by
\(B_\chi\) times the variation of the zero extension of
\(\psi_j\) on the positive half-line.
It also vanishes for \(x>\Lambda\).
Thus \(e_j\) lies in \(L^1(dx/x)\cap L^2(dx/x)\);
near zero it is \(O(\sqrt x)\).

For \(\operatorname{Re}s>1\), absolute convergence gives
\[
 \int_0^\infty\left(\sum_{n\geq1}\chi(n)\psi_j(nx)\right)
                      x^{s-1}\,dx
       =L(s,\chi)\int_0^\infty\psi_j(x)x^{s-1}\,dx.
\]
The summation-by-parts bound continues the left side
holomorphically to \(\operatorname{Re}s>0\); the right side
has the same continuation.  At \(s=1/2+i\tau\) this is the
Fourier transform of \(e_j(e^t)\).
Plancherel proves (TB.1).  The finite-rank projection formula
\(Q_\Lambda=EG^{-1}E^*\) from AC-T11 gives (TB.2).
The integral of \(K_{\chi,\Lambda}\) divided by \(2\pi\) is
\(\operatorname{tr}(G^{-1}G)=N\), which proves the claimed
absolute convergence.

If \(\psi(x)=q^{-1/4}\phi(x/\sqrt q)\), its Mellin
transform at \(1/2+i\tau\) is \(q^{i\tau/2}\) times that of
\(\phi\).  The common phases cancel in both matrix formulas.
Finally suppose the two stated weights were proportional
with constant \(c>0\).  Analytic continuation from the line
\(\operatorname{Re}s=1/2\) would imply
\[
 L(s,\chi)L(1-s,\overline\chi)
       =c\,\zeta(s)\zeta(1-s).
\]
The left side is entire.  The right side has a pole at \(s=1\),
because \(\zeta(0)=-1/2\).  This is impossible.
\end{proof}

\begin{proposition}[AC-T22: exact tail and principal-angle errors]
\label{AC-T22}
Let
\[
 H=E^*(1-S_\Lambda)E,\qquad
 A_\Lambda=G^{-1/2}HG^{-1/2}.
\]
Then \(0\leq A_\Lambda<1\).
If \(Q_\Lambda^0\) denotes the orthogonal projection onto
\(\operatorname{span}\{S_\Lambda e_j\}\), then
\[
 \|Q_\Lambda-Q_\Lambda^0\|_2^2
       =2\operatorname{tr}A_\Lambda,\qquad
 \|Q_\Lambda-Q_\Lambda^0\|_1
       =2\operatorname{tr}\sqrt{A_\Lambda}
       \leq2\sqrt{N\operatorname{tr}A_\Lambda}.
 \tag{TB.3}
\]
In particular
\[
 \left|\operatorname{Tr}\bigl((Q_\Lambda-Q_\Lambda^0)V(h)\bigr)\right|
       \leq 2\|h\|_{L^1(dx/x)}
                     \operatorname{tr}\sqrt{A_\Lambda}.
 \tag{TB.4}
\]
The tail matrix has the exact conductor-sensitive expression
\[
 H_{ij}=\int_{\Lambda/q}^\infty
   \overline{E_{\overline\chi}(\widehat\psi_i)(y)}
             E_{\overline\chi}(\widehat\psi_j)(y)\,\frac{dy}{y}.
 \tag{TB.5}
\]
For piecewise smooth cutoff vectors the theta sums on the
right can equivalently be understood through the Poisson
identity almost everywhere.
\end{proposition}
\begin{proof}
Every \(e_j\) vanishes above \(\Lambda\), so \(H\) is
precisely the Gram matrix of the tails below \(\Lambda^{-1}\).
Finite Fourier transformation of \(\chi\), followed by ordinary
Poisson summation, gives
\[
 E_\chi\psi(x)=\omega_\chi
       E_{\overline\chi}\widehat\psi(1/(qx)),
 \qquad |\omega_\chi|=1.
 \tag{TB.6}
\]
The phase is the normalized Gauss sum, with the harmless
sign determined by the Fourier convention.  There are no
constant terms because \(\chi\) is nonprincipal.
For compactly supported piecewise smooth functions this is
the ordinary Poisson identity for a periodized function of
bounded variation, away from its countable endpoint set.
Changing \(y=1/(qx)\) in the tail integral proves (TB.5).

The vectors \(S_\Lambda e_j\) are linearly independent.
Indeed on
\((\max(\Lambda^{-1},\Lambda/2),\Lambda)\) the theta sum
has only its \(n=1\) term.  If a linear combination of the
truncated images vanishes there, the same combination of
prolate functions vanishes on that open interval.
These functions are analytic inside \((-\Lambda,\Lambda)\),
so the combination is identically zero.  Their original
linear independence gives the assertion.
Thus \(G-H\) is positive definite and \(A_\Lambda<1\).

Put \(Z=EG^{-1/2}\), an isometry from \(\mathbb C^N\).
Then \(Z^*S_\Lambda Z=1-A_\Lambda\) and
\[
 Q_\Lambda^0=S_\Lambda Z(1-A_\Lambda)^{-1}Z^*S_\Lambda.
\]
Diagonalize \(A_\Lambda\) with eigenvalues \(\alpha_j\).
The pairs consisting of \(Z u_j\) and
\(S_\Lambda Zu_j/\sqrt{1-\alpha_j}\) are orthogonal for
different \(j\), with mutual inner product
\(\sqrt{1-\alpha_j}\) in the \(j\)-th pair.
On their two-dimensional span the difference of the
rank-one projections has eigenvalues
\(\pm\sqrt{\alpha_j}\).  Zero angles contribute nothing.
Summation proves both norm identities in (TB.3), and
Cauchy--Schwarz proves its final inequality.
Finally \(\|V(h)\|\leq\|h\|_1\) gives (TB.4).
\end{proof}

\begin{proposition}[AC-T23: a computable sufficient trace gate]
\label{AC-T23}
Write \(L=\log\Lambda\) and let
\(k^0_{\chi,\Lambda}(\tau)\) be the Fourier diagonal density
of \(Q_\Lambda^0\), defined from any orthonormal basis of its
range.  Then
\[
 0\leq k^0_{\chi,\Lambda}(\tau)\leq2L,\qquad
 \operatorname{Tr}\bigl((S_\Lambda-Q_\Lambda^0)V(h)\bigr)
 =\frac1{2\pi}\int_{\mathbb R}
       \widehat h_+(\tau)\bigl(2L-k^0_{\chi,\Lambda}(\tau)\bigr)
       \,d\tau.
 \tag{TB.7}
\]
Consequently the following two estimates suffice for the
nonprincipal-character assertion T-CUT:
\[
 \operatorname{tr}\sqrt{A_\Lambda}=o(1),\qquad
 \frac{(2L-k^0_{\chi,\Lambda}(\tau))\,d\tau}{2\pi}
  \ \longrightarrow\ \sum_\gamma m_{\chi,\gamma}\delta_\gamma
 \tag{TB.8}
\]
where the second convergence is tested against
\(\widehat h_+\) for every \(h\in C_c^\infty(\mathbb R_{>0})\).
The second estimate includes control at high frequency,
not just vague convergence on compact sets.
The first estimate is a sufficient trace-norm condition,
not a necessary or equivalent reformulation of T-CUT.
No decay of this quantity is asserted at C99's full
threshold.  The transition of the scalar prolate eigenvalues
described after C99, VIII.(32), makes it particularly important
not to replace that threshold silently by a smaller family.
AC-T25--AC-T26 below give a comparison after testing the
traces and retain the original family.

A direct variation estimate gives the fully explicit bound
\[
 \operatorname{tr}A_\Lambda
 \leq\frac{B_\chi^2}{\Lambda\,\lambda_{\min}(G)}
        \sum_{j=1}^N
          \operatorname{Var}_{(0,\infty)}(\psi_j)^2 .
 \tag{TB.9}
\]
This estimate alone does not establish (TB.8) for the
growing prolate family of size \(N\asymp\Lambda^2/q\).
\end{proposition}
\begin{proof}
In logarithmic coordinates let \(z_1,\ldots,z_N\) be an
orthonormal basis of the range of \(Q_\Lambda^0\),
supported in \([-L,L]\).
Then
\[
 k^0_{\chi,\Lambda}(\tau)
       =\sum_{j=1}^N\left|
                  \int_{-L}^Lz_j(t)e^{i\tau t}\,dt\right|^2.
 \]
Bessel's inequality against the vector
\(t\mapsto e^{-i\tau t}\) in \(L^2[-L,L]\) gives the upper
bound \(2L\).  Fourier inversion and the finite-rank trace
give (TB.7); the interval convolution trace is
\(2Lh(1)\), as in AC-T19.
Combining (TB.7) and (TB.4) proves sufficiency of (TB.8).

For (TB.9), summation by parts as in AC-T21 gives
\[
 |e_j(x)|\leq B_\chi\sqrt{x}\,
                     \operatorname{Var}_{(0,\infty)}(\psi_j).
 \]
Integration on \(0<x<\Lambda^{-1}\) bounds \(H_{jj}\)
by \(B_\chi^2\operatorname{Var}(\psi_j)^2/\Lambda\).
Since \(G^{-1}\leq\lambda_{\min}(G)^{-1}\), tracing gives
(TB.9).  The estimate (TB.9) by itself supplies no sufficient control
through the prolate threshold.  AC-T27--AC-T29 below obtain
quantitative inverse-Gram and smoothed-tail bounds by different
arguments; the zero-measure limit remains (TB.53).
\end{proof}

The second limit in (TB.8) is the continuous counterpart of the
finite-zero computation in C99, Lemma VIII.3, equations
(28)--(29).  In the function-field proof only finitely many
zeros occur in the fundamental frequency domain, and a dimension
count identifies the image exactly.
For a number-field character there are infinitely many zeros.
Under \(A\), compactly supported logarithmic functions cannot vanish at
all of them unless zero, by AC-T13's Jensen argument.
The finite image is therefore not characterized by imposing
all those exact vanishing conditions.
Equations (TB.1), (TB.5), and (TB.8) identify the additional
arithmetic estimates needed after the scalar prolate rescaling.
They do not assert that these estimates are false, or that they
are logically independent of \(A\).

\subsection{An exact closure consequence and why its cutoff is different}

\begin{corollary}[AC-T24: one-sided closures under \(A\)]
\label{AC-T24}
In the preceding nonprincipal character sector, \(A\) gives
\[
 \overline{E_\chi(C_c^\infty(-R,R)^\epsilon)}
       =L^2((0,R),dx/x),\qquad R>0.
 \tag{TB.10}
\]
Here the superscript fixes parity.  It follows that
\[
 \overline{E_\chi(C_c^\infty(-\Lambda,\Lambda)^\epsilon)}
 \ \cap\
 \overline{E_\chi\bigl(\mathcal F^{-1}
              C_c^\infty(-\Lambda/q,\Lambda/q)^\epsilon\bigr)}
       =L^2([\Lambda^{-1},\Lambda],dx/x).
 \tag{TB.11}
\]
\end{corollary}
\begin{proof}
Burnol, \emph{An adelic causality problem related to abelian
\(L\)-functions}, arXiv:\path{math/0001013v3},
Theorem 3.2 and its number-field proof on pp.~13--15,
identifies this one-sided closure with the Hardy subspace
multiplied by the Blaschke product of the zeros to the right
of the critical line.  Burnol's multiplier \(V\) is identically
one on every nonprincipal character sector (display on
printed p.~13), so no extra operator modifies this closure.
The fixed finite character vector and
arbitrary smooth archimedean factors used there suffice.
Under \(A\) that product is one, proving (TB.10) at \(R=1\);
dilation gives arbitrary \(R\).
Apply (TB.10) to \(\overline\chi\) and use (TB.6) to obtain
the second closed space in (TB.11) as
\(L^2((\Lambda^{-1},\infty),dx/x)\).
Intersecting the two support spaces proves (TB.11).
\end{proof}

The completed intersection in (TB.11) is the whole norm window.
Its projection is \(S_\Lambda\), whose subtracted trace defect
is identically zero.  The prolate-image projection is the
finite-rank \(Q_\Lambda\); its proposed defect contains the
critical-zero measure.  Thus replacing the finite approximation
by the intersection of the completed one-sided ranges would
discard the spectral information that T-CUT seeks to recover.
The exact closure result is useful and is a consequence of \(A\);
the rate and trace comparison in (TB.8) remain separate work.

\subsection{A uniform smoothed theta estimate}

\begin{lemma}[AC-T25: smoothing a half-line theta image]
\label{AC-T25}
Let \(\chi\) be primitive and nonprincipal.  Define the closed
theta operator \(\mathcal E_\chi:L^2(\mathbb R_{>0},dx)
\longrightarrow L^2(\mathbb R_{>0},dx/x)\) by Mellin
multiplication by \(L(1/2+i\tau,\chi)\), with its maximal
multiplier domain.  On compactly supported functions of bounded
variation this is the theta sum used in AC-T21.
For every \(g\in C_c^\infty(\mathbb R_{>0})\) there is a
constant \(C_{\chi,g}\), independent of \(r>0\), such that
\[
 \|V(g)1_{(r,\infty)}\mathcal E_\chi\psi\|_{L^2(dx/x)}
       \leq C_{\chi,g}\|\psi\|_{L^2(dx)}.
 \tag{TB.12}
\]
The operator on its original multiplier domain therefore
extends boundedly to all of \(L^2(dx)\).

More generally (TB.12) holds for \(g\in L^1(dx/x)\cap
L^2(dx/x)\) whenever, for \(\theta=5/16\) and
\(m=\widehat g_+\),
\[
 \|m(\tau)L(1/2+i\tau,\chi)\|_\infty
       +\|m(\tau)(1+|\tau|)^\theta\|_{L^2(d\tau)}<\infty.
 \tag{TB.13}
\]
One can take \(C_{\chi,g}\) bounded by this sum with a
constant depending only on \(\chi\).

Normalize the existing prolate vectors to be orthonormal on
the positive half-line, without changing their span or their
indices.  Let \(\nu_j\) be their concentration eigenvalues
for time radius \(\Lambda\) and frequency radius
\(R=\Lambda/q\), and put
\(D_\Lambda=\operatorname{diag}(1-\nu_j)\).
Then for every such smooth \(g\),
\[
 E^*(1-S_\Lambda)V(g)^*V(g)(1-S_\Lambda)E
                \leq C_{\overline\chi,g^\vee}^{\,2}D_\Lambda,
 \qquad g^\vee(a)=g(a^{-1}),
 \tag{TB.14}
\]
and, with \(Z=EG^{-1/2}\),
\[
 \|V(g)(1-S_\Lambda)Z\|_2^2
       \leq C_{\overline\chi,g^\vee}^{\,2}
                       \operatorname{tr}(G^{-1}D_\Lambda).
 \tag{TB.15}
\]
The norm in (TB.15) is the Hilbert--Schmidt norm of a map
from \(\mathbb C^N\).  These estimates are unconditional.
\end{lemma}
\begin{proof}
Write \(L_\chi(\tau)=L(1/2+i\tau,\chi)\).
We first record the elementary fixed-character convexity
bound in a form sufficient for this argument:
\[
 |L_\chi(\tau)|+|L_\chi'(\tau)|
                      \leq C_\chi(1+|\tau|)^{5/16}.
 \tag{TB.16}
\]
Here the prime denotes the derivative with respect to \(\tau\).
Indeed the absolutely convergent Dirichlet series is bounded
on \(\operatorname{Re}s=17/16\).  The functional equation
with gamma factor \(\Gamma((s+\epsilon)/2)\), together with
Stirling's formula, bounds it on
\(\operatorname{Re}s=-1/16\) by
\(C_\chi(1+|\operatorname{Im}s|)^{9/16}\).
The polynomial Phragmen--Lindelof convexity theorem for this
finite-order entire function bounds it throughout
\(7/16\leq\operatorname{Re}s\leq9/16\) by the power
\(5/16\).  Cauchy's formula on a circle of radius \(1/32\)
about \(1/2+i\tau\) gives the derivative bound.  Bounded
values of \(\tau\) are absorbed in the constant.  This use
of the functional equation is equivalently obtained from
the character Poisson identity (TB.6); it uses no zero
location hypothesis.

The unitary map \(\psi\mapsto e^{t/2}\psi(e^t)\), followed
by the unitary logarithmic Fourier transform, identifies
\(\mathcal E_\chi\) with multiplication by \(L_\chi\).
In the target logarithmic variable let
\(P_a=1_{(a,\infty)}\), where \(a=\log r\).
At \(a=0\), its Fourier representation is the Hardy
projection with kernel
\(\tfrac12\delta(\tau-\sigma)+
\tfrac{i}{2\pi}\operatorname{pv}(\tau-\sigma)^{-1}\).
Consequently, initially on compactly supported smooth Fourier
vectors,
\[
 M_mP_0M_{L_\chi}
       =M_{mL_\chi}P_0+K_m,
 \qquad
 K_m(\tau,\sigma)=\frac{i}{2\pi}m(\tau)
       \frac{L_\chi(\sigma)-L_\chi(\tau)}{\tau-\sigma}.
 \tag{TB.17}
\]
The delta contribution cancels in the commutator.
For \(|\sigma-\tau|\leq1\), (TB.16) bounds the divided
difference by \(C_\chi(1+|\tau|)^\theta\).
For \(|\sigma-\tau|>1\), its square is bounded by
\[
 C_\chi\left((1+|\tau|)^{2\theta}
                  |\sigma-\tau|^{-2}
             +|\sigma-\tau|^{2\theta-2}\right).
\]
Both terms are integrable in \(\sigma\), because
\(2\theta<1\).  Hence
\[
 \|K_m\|_2^2
       \leq C_\chi\int_{\mathbb R}
                    |m(\tau)|^2(1+|\tau|)^{2\theta}\,d\tau.
\]
The other term in (TB.17) has norm at most
\(\|mL_\chi\|_\infty\).  For general \(a\), the Hardy
projection is conjugated by multiplication by a unit phase;
that phase commutes with \(M_m\) and \(M_{L_\chi}\).
The bounds are therefore independent of \(a\).
Smooth compact Fourier vectors are a graph core for
\(M_{L_\chi}\).  Approximation in that graph norm shows that
the bounded extension agrees with the stated operator on
its full original domain.  This proves (TB.12)--(TB.13).

Set \(\eta_j=(1-1_{[-R,R]})\widehat\psi_j\), retaining
the prescribed parity.  The concentration eigenvalue
equation and Plancherel give, in the positive-half-line
normalization,
\[
 \langle\eta_i,\eta_j\rangle_{L^2(dx)}
                  =(1-\nu_j)\delta_{ij}.
 \tag{TB.18}
\]
All theta identities needed for these vectors hold in the
closed multiplier domains.  To see this without assigning
a pointwise meaning to a conditionally convergent tail sum,
approximate each compact prolate cutoff by smooth compact
functions of the same parity, with uniformly bounded
variation and a common compact support.  Their original
theta images converge in \(L^2(dx/x)\), by the Abel bound
in AC-T21 and dominated convergence.  Poisson summation for
the approximants and closedness of
\(\mathcal E_{\overline\chi}\) then show that
\(\widehat\psi_j\) belongs to its domain and satisfies
(TB.6).  Its frequency truncation is compactly supported
and of bounded variation, so it too belongs to that domain;
therefore \(\eta_j\) does as well.

Let \((J_qf)(x)=f(1/(qx))\), a unitary operator on
\(L^2(dx/x)\).  Since the theta image of
\(1_{[-R,R]}\widehat\psi_j\) vanishes for \(y>R\),
the exact tail identity becomes
\[
 (1-S_\Lambda)Ec
    =\omega_\chi J_q1_{(R,\infty)}
                  \mathcal E_{\overline\chi}\sum_jc_j\eta_j.
 \tag{TB.19}
\]
The original theta image is zero above \(\Lambda\), so
there is no second tail on the left.  Also
\(V(g)J_q=J_qV(g^\vee)\).
Apply (TB.12) and then (TB.18) to (TB.19).
The result is (TB.14) for every vector \(c\), hence the
matrix inequality.  Conjugation by \(G^{-1/2}\) and tracing
give (TB.15).  The same argument applies to the larger
test class in (TB.13) whenever its reflected test satisfies
that condition for \(\overline\chi\).
\end{proof}

\begin{proposition}[AC-T26: exact comparison after testing the trace]
\label{AC-T26}
Keep the full prolate family specified above.  Put
\(Z_0=S_\Lambda Z\), \(Z_1=(1-S_\Lambda)Z\), and
\(A=A_\Lambda=Z_1^*Z_1\).
For \(g\in L^1(dx/x)\cap L^2(dx/x)\), write
\[
 C_g=V(g)^*Z_1,\qquad D_g=V(g)^*Z_0,\qquad
 \varepsilon_g=\|C_g\|_2^2,\qquad
 r_g=\|D_gA^{1/2}(1-A)^{-1/2}\|_2^2.
\]
Then the tested comparison is the exact identity
\[
 \Delta_\Lambda(g):=
 \operatorname{Tr}\bigl((Q_\Lambda-Q_\Lambda^0)V(g*g^*)\bigr)
 =\varepsilon_g+2\operatorname{Re}\operatorname{tr}(C_g^*D_g)-r_g.
 \tag{TB.20}
\]
In particular
\[
 |\Delta_\Lambda(g)|
 \leq\varepsilon_g+
         2\sqrt{2\log\Lambda}\,\|g\|_2\sqrt{\varepsilon_g}+r_g,
 \qquad
 r_g\leq2\log\Lambda\,\|g\|_2^2
                      \frac{\|A\|}{1-\|A\|}.
 \tag{TB.21}
\]
For smooth compact tests, AC-T25 gives
\[
 \varepsilon_g\leq
 C_{\overline\chi,(g^*)^\vee}^{\,2}
                         \operatorname{tr}(G^{-1}D_\Lambda).
 \tag{TB.22}
\]
All terms in (TB.20) are computable from the original Gram
matrices.  More explicitly, define
\[
 \begin{split}
 T_g&=E^*(1-S_\Lambda)V(g*g^*)(1-S_\Lambda)E,\\
 X_g&=E^*(1-S_\Lambda)V(g*g^*)S_\Lambda E,\\
 M_g^0&=E^*S_\Lambda V(g*g^*)S_\Lambda E.
 \end{split}
\]
Then
\[
 \Delta_\Lambda(g)=
 \operatorname{tr}(G^{-1}T_g)
 +2\operatorname{Re}\operatorname{tr}(G^{-1}X_g)
 -\operatorname{tr}\left(((G-H)^{-1}-G^{-1})M_g^0\right).
 \tag{TB.23}
\]
For example, the sufficient estimates
\[
 \varepsilon_g=o((\log\Lambda)^{-1}),\qquad r_g=o(1)
 \tag{TB.24}
\]
would give \(\Delta_\Lambda(g)\to0\).
The first follows from the stronger inverse-Gram estimate
\(\operatorname{tr}(G^{-1}D_\Lambda)=o((\log\Lambda)^{-1})\).
Neither that estimate nor (TB.24) is asserted here at the
full threshold.  In particular, the raw scalar prolate
transition estimate by itself gives no such bound for the
arithmetic inverse Gram matrix.
AC-T29 below proves the first estimate in (TB.24) by a
sharper argument.  AC-T30 then avoids the term \(r_g\) by
using the compressed positive contraction instead of its
range projection.

An exact tested trace condition in this range-projection
comparison can be written
without requiring separate decay of its three error terms:
\[
 \frac1{2\pi}\int_{\mathbb R}|\widehat g_+(\tau)|^2
                 (2\log\Lambda-k^0_{\chi,\Lambda}(\tau))\,d\tau
       -\Delta_\Lambda(g)
 \longrightarrow
       \sum_\gamma m_{\chi,\gamma}|\widehat g_+(\gamma)|^2.
 \tag{TB.25}
\]
For \(g\in C_c^\infty\), this is exactly T-CUT on the
positive convolution test \(g*g^*\), not an additional
claimed consequence of \(A\).
\end{proposition}
\begin{proof}
The two projection formulas are
\(Q_\Lambda=ZZ^*\) and
\(Q_\Lambda^0=Z_0(1-A)^{-1}Z_0^*\).
Since \(V(g*g^*)=V(g)V(g)^*\), their traces are
\(\|C_g+D_g\|_2^2\) and
\(\operatorname{tr}((1-A)^{-1}D_g^*D_g)\).
Subtracting gives (TB.20), because
\((1-A)^{-1}-1=A(1-A)^{-1}\).

The operator inequality \(Z_0Z_0^*=S_\Lambda Q_\Lambda
S_\Lambda\leq S_\Lambda\) implies
\[
 \|D_g\|_2^2
 \leq\|V(g)^*S_\Lambda\|_2^2
              =2\log\Lambda\,\|g\|_{L^2(dx/x)}^2.
\]
The last equality follows directly by integrating the squared
logarithmic convolution kernel over an interval of length
\(2\log\Lambda\).  Hilbert--Schmidt Cauchy--Schwarz bounds
the cross term in (TB.20); the norm of
\(A^{1/2}(1-A)^{-1/2}\) gives the other inequality in
(TB.21).  Applying (TB.15) to \(g^*\) proves (TB.22).
Conjugating the three finite matrices by \(G^{-1/2}\)
gives (TB.23), using
\[
 (G-H)^{-1}-G^{-1}
 =G^{-1/2}\bigl((1-A)^{-1}-1\bigr)G^{-1/2}.
\]
Equations (TB.21)--(TB.22) prove the stated sufficient
conditions.  Finally (TB.7), with
\(\widehat{g*g^*}_+=|\widehat g_+|^2\), gives the exact
left side of (TB.25) as
\(2\log\Lambda\,\|g\|_2^2-
\operatorname{Tr}(Q_\Lambda V(g*g^*))\).
This proves the asserted equivalence on these tests.
\end{proof}

\begin{remark}[All smooth compact trace tests]
\label{rem:AC-T26alltests}
The comparison \(\Delta_\Lambda(g)\to0\) for quadratic
tests can cover every original \(h\in C_c^\infty\) without
assuming an unproved uniform distribution bound.  It suffices
to establish it on the linear space generated by smooth
compact tests and the single logarithmic kernel
\(k(e^t)=\tfrac12e^{-|t|}\).
Indeed \(\widehat k_+(\tau)=(1+\tau^2)^{-1}\), so \(k\)
and its sums with compact smooth tests satisfy (TB.13), as
do their reflections and involutions.  If
\(H(t)=h(e^t)\), set
\(j(e^t)=(1-\partial_t^2)\overline{H(-t)}\).
Fourier transformation gives \(h=k*j^*\).
Polarization of the Hermitian sesquilinear form
\[
 (g_1,g_2)\longmapsto
 \operatorname{Tr}\bigl((Q_\Lambda-Q_\Lambda^0)
                            V(g_1)V(g_2)^*\bigr)
\]
then gives convergence for this \(h\).
The hypotheses in this remark concern the tested comparison,
not trace-norm convergence of the projections.
\end{remark}

\subsection{Bounds for the actual arithmetic inverse Gram matrix}

\begin{proposition}[AC-T27: an unconditional explicit Gram lower bound]
\label{AC-T27}
Let \(c=\Lambda R=\Lambda^2/q\geq1\), and let
\(\nu=\min_j\nu_j>0\) for the original selected prolate family.
With the positive-half-line orthonormal normalization of
AC-T25, its arithmetic Gram matrix satisfies
\[
       G\geq\frac{\nu}{64c}\,1,\qquad
       \|G^{-1}\|\leq\frac{64c}{\nu}.
 \tag{TB.26}
\]
This bound requires neither \(A\) nor any asymptotic assertion
about the concentration eigenvalues.
\end{proposition}
\begin{proof}
Let \(P\) be the time cutoff on \([-\Lambda,\Lambda]\),
\(Q\) the frequency cutoff on \([-R,R]\), and \(\Pi\) the
orthogonal projection onto the selected prolate functions in
the full real-line normalization.  Their eigenvalue equations
give \(\Pi\leq\nu^{-1}PQP\).  The continuous diagonal kernel
of \(PQP\) in the interior of the time interval is \(2R\).
Thus every vector \(\psi\) of the selected parity satisfies
\[
 |\psi(x)|^2\leq\frac{4R}{\nu}
                     \|\psi\|_{L^2(\mathbb R_{>0},dx)}^2
       \quad(0<x<\Lambda).
 \tag{TB.27}
\]
The factor two comes from passing from the full-line norm to
the positive-half-line norm.  For
\(a=\nu/(8R)\), we have \(a<\Lambda\) and consequently
\[
 \|\psi\|_{L^2(a,\Lambda)}^2
                     \geq\tfrac12\|\psi\|_{L^2(0,\Lambda)}^2.
 \tag{TB.28}
\]

Write \(e=E_\chi\psi\).  Finite Dirichlet convolution gives,
pointwise away from the finitely many endpoints in each sum,
\[
 \sqrt{x}\psi(x)
       =\sum_{n\leq\Lambda/x}
                     \frac{\mu(n)\chi(n)}{\sqrt n}\,e(nx).
 \tag{TB.29}
\]
Indeed expansion of \(e(nx)\) makes the coefficient of
\(\psi(kx)\) equal to
\(\chi(k)\sum_{n\mid k}\mu(n)\), which is zero unless \(k=1\).
For \(x\geq a\) the sum has at most \(\lfloor\Lambda/a\rfloor\)
terms.  The multiplicative translation \(e(x)\mapsto e(nx)\)
is unitary on \(L^2(dx/x)\).  Minkowski's inequality and
\(\sum_{n\leq M}n^{-1/2}\leq2\sqrt M\) therefore give
\[
 \|\psi\|_{L^2(a,\Lambda)}
       \leq2\sqrt{\Lambda/a}\,\|e\|_{L^2(dx/x)}.
 \tag{TB.30}
\]
Combining (TB.28)--(TB.30) yields
\(\|e\|^2\geq\nu\|\psi\|^2/(64\Lambda R)\).
Because the original \(\psi_j\) are orthonormal on the
positive half-line, this is exactly (TB.26).
\end{proof}

\begin{theorem}[AC-T28: subpower inverse-Gram control from \(A\)]
\label{AC-T28}
Assume \(A\), and keep the fixed primitive nonprincipal
character \(\chi\).  For every \(\kappa>0\) there is a constant
\(C_{\chi,\kappa}>0\), independent of \(\Lambda\) and of the
selected prolate indices, such that
\[
       G\geq C_{\chi,\kappa}\left(\frac{\nu}{c}\right)^\kappa 1,
       \qquad
       \|G^{-1}\|\leq
              C_{\chi,\kappa}^{-1}\left(\frac c\nu\right)^\kappa
       \quad(c\geq1).
 \tag{TB.31}
\]
In particular, on any original cutoff sequence for which
the scalar concentration eigenvalue \(\nu\) stays bounded
below, the arithmetic inverse Gram matrix grows more slowly
than any prescribed positive power of \(c\).
This conclusion concerns \(G\); it does not replace the
truncated-image matrix \(G-H\) in (TB.23).
\end{theorem}
\begin{proof}
We first prove the reciprocal bound actually needed.  For
each fixed \(0<\delta<1/32\) and each \(\eta>0\), \(A\) gives
\[
 \left|L(1/2+\delta+i\tau,\chi)^{-1}\right|
                   \leq C_{\chi,\delta,\eta}(1+|\tau|)^\eta.
 \tag{TB.32}
\]
Indeed the nonprincipal \(L\)-function is entire and has no
zeros in \(\operatorname{Re}s>1/2\).  Its Euler logarithm
therefore extends to an analytic logarithm on that simply
connected half-plane.  Center disks at \(2+i\tau\), with radii
\[
 r_0=\tfrac12,\qquad r_1=\tfrac32-\delta,\qquad
 r_2=\tfrac32-\tfrac\delta2,\qquad
 r_3=\tfrac32-\tfrac\delta4.
\]
The logarithm is uniformly bounded on the disk of radius
\(r_0\), by its absolutely convergent Euler series.
The ordinary polynomial vertical-strip bound for \(L\),
as used in (TB.16), bounds the real part of the logarithm
on the largest disk by \(C_\chi\log(2+|\tau|)\).
Borel--Caratheodory bounds the modulus of the logarithm on
the disk of radius \(r_2\) by
\(C_{\chi,\delta}\log(2+|\tau|)\).
Hadamard's three-circles theorem, between \(r_0,r_1,r_2\),
then bounds it at \(1/2+\delta+i\tau\) by
\[
 C_{\chi,\delta}\bigl(\log(2+|\tau|)\bigr)^{\beta_\delta},
 \qquad
 \beta_\delta=\frac{\log(r_1/r_0)}{\log(r_2/r_0)}<1.
\]
Exponentiation gives (TB.32), since every power of a
logarithm with exponent strictly below one is
\(o(\log(2+|\tau|))\).  This proof uses exactly the zero-free
half-plane supplied by \(A\).

We next record a uniform fractional Sobolev estimate.
For \(0<\alpha<1/2\), a function
\(v\in H^1((-\infty,b))\), extended by zero above \(b\),
satisfies
\[
 \|v\|_{H^\alpha(\mathbb R)}
 \leq C_\alpha\|v\|_2^{1-\alpha}
                         (\|v\|_2+\|v'\|_2)^\alpha.
 \tag{TB.33}
\]
Here and below \(H^\alpha\) uses Fourier weight
\((1+\tau^2)^\alpha\) in the squared norm.
For completeness, integration by parts gives
\[
 i\tau\widehat v_+(\tau)
       =v(b)e^{i\tau b}
                    -\widehat{1_{(-\infty,b)}v'}_+(\tau),
 \qquad |v(b)|^2\leq2\|v\|_2\|v'\|_2.
\]
Split the Fourier integral at
\(K=1+\|v'\|_2/\|v\|_2\).
The part \(|\tau|\leq K\) is at most
\(C K^{2\alpha}\|v\|_2^2\).
On the complement, the derivative term is at most
\(C K^{2\alpha-2}\|v'\|_2^2\), and the endpoint term is at
most \(C_\alpha K^{2\alpha-1}\|v\|_2\|v'\|_2\).
The last integral is finite precisely because \(\alpha<1/2\).
These bounds prove (TB.33); the zero function is immediate.

For a selected vector \(\psi=\sum_jb_j\psi_j\), put
\(\phi=\sum_jb_j\nu_j^{-1}\psi_j\).
The prolate equation implies \(\psi=Q\phi\) inside the time
interval, while
\(\|\phi\|_2\leq\nu^{-1}\|\psi\|_2\).
The band limitation of \(Q\phi\) therefore gives
\[
       \|\psi'\|_{L^2(0,\Lambda)}
             \leq\frac{2\pi R}{\nu}\|\psi\|_{L^2(0,\Lambda)}.
 \tag{TB.34}
\]
The same inequality on the positive half-line follows from
parity and full-line Plancherel.

Write \(e=E_\chi\psi\), \(b=\log\Lambda\), and define
\[
 v_\delta(t)=\Lambda^{-\delta}
               e^{(1/2+\delta)t}\psi(e^t),\qquad
 w_\delta(t)=\Lambda^{-\delta}e^{\delta t}e(e^t).
\]
Both are zero for \(t>b\).  Mellin transformation gives
\[
 \widehat w_{\delta,+}(\tau)
        =L(1/2+\delta+i\tau,\chi)
                             \widehat v_{\delta,+}(\tau).
 \tag{TB.35}
\]
The analytic-continuation argument in AC-T21 proves this
identity on this vertical line for the present vectors.
We have
\[
 \|w_\delta\|_2\leq\|e\|_2,\qquad
 \|v_\delta\|_2\leq\|\psi\|_2,\qquad
 \|v_\delta'\|_{L^2(-\infty,b)}
                    \leq C\left(1+\frac c\nu\right)\|\psi\|_2.
 \tag{TB.36}
\]
The last bound follows by differentiating and using (TB.34).
On the other hand, (TB.28), with \(a=\nu/(8R)\), gives
\[
 \|v_\delta\|_2\geq
          2^{-1/2}\left(\frac{\nu}{8c}\right)^\delta\|\psi\|_2.
 \tag{TB.37}
\]

Fix \(s=1/16\) and \(\alpha=3/8=s+5/16\).
The unconditional convexity bound used in (TB.16) holds
on \(\operatorname{Re}s=1/2+\delta\) with exponent \(5/16\).
Combining (TB.33)--(TB.36) thus gives
\[
 \|w_\delta\|_{H^s}
            \leq C_{\chi,\delta}
                       \left(1+\frac c\nu\right)^\alpha\|\psi\|_2.
 \tag{TB.38}
\]
Choose \(0<\eta<s\) and put \(d=\eta/s\).
By (TB.32), (TB.35), and Fourier Holder interpolation,
\[
 \begin{split}
 \|v_\delta\|_2
   &\leq C_{\chi,\delta,\eta}\|w_\delta\|_{H^\eta}\\
   &\leq C_{\chi,\delta,\eta}
        \|e\|_2^{\,1-d}
        \left(1+\frac c\nu\right)^{\alpha d}\|\psi\|_2^d .
 \end{split}
 \tag{TB.39}
\]
Together with (TB.37), this proves
\[
 \|e\|_2^2\geq
 C_{\chi,\delta,d}
 \left(\frac{\nu}{c}\right)^{(2\delta+2\alpha d)/(1-d)}
                                         \|\psi\|_2^2,
 \tag{TB.40}
\]
where \(c/\nu\geq1\) absorbs the harmless constants.
Both \(\delta\) and \(d\) can be chosen arbitrarily small
and positive.  Choose them so that the exponent in
(TB.40) is at most the prescribed \(\kappa\).
Since \(\nu/c\leq1\), this gives (TB.31) for every vector
of the selected family.
\end{proof}


\subsection{Decay of the tested tail at the original prolate threshold}

\begin{theorem}[AC-T29: the actual smoothed theta tail tends to zero]
\label{AC-T29}
Let
\[
 k(e^t)=\tfrac12e^{-|t|},\qquad
 \mathscr G=C_c^\infty(\mathbb R_{>0})+\mathbb C k .
\]
For \(g\in\mathscr G\), the original positive-half-line
orthonormal prolate family satisfies the unconditional bound
\[
 \sum_{j=1}^N
       \|V(g)(1-S_\Lambda)E_\chi\psi_j\|_2^2
                       \leq\frac{C_{\chi,g}}{\nu\Lambda^2}.
 \tag{TB.41}
\]
Under \(A\), for every \(\kappa>0\) this gives
\[
 \|V(g)(1-S_\Lambda)EG^{-1/2}\|_2^2
       \leq\frac{C_{\chi,g,\kappa}}{\nu\Lambda^2}
                                  \left(\frac c\nu\right)^\kappa.
 \tag{TB.42}
\]
In particular, for the full parity-index cutoff
\(n\leq4\Lambda^2/q\) used in the conductor rescaling of C99,
\[
 \|V(g)(1-S_\Lambda)EG^{-1/2}\|_2^2
                 =O_{\chi,g,\kappa}(\Lambda^{-2+2\kappa}).
 \tag{TB.43}
\]
No prolate transition indices are deleted in (TB.43).
\end{theorem}
\begin{proof}
On the raw positive-half-line space with the prescribed parity,
write \(P=1_{[-\Lambda,\Lambda]}\) and
\(Q=1_{[-R,R]}(\mathcal F)\), with \(R=\Lambda/q\).
Let
\[
       T=V(g)1_{(0,\Lambda^{-1})}\mathcal E_\chi .
\]
This has a bounded extension by AC-T25: subtract the
upper-half-line operator from \(V(g)\mathcal E_\chi\),
whose Mellin multiplier \(mL_\chi\) is bounded.
The same argument applies to \(k\), since its multiplier
is \((1+\tau^2)^{-1}\), and hence to \(\mathscr G\).

The exact character Poisson identity implies \(TQ=0\).
Initially take a Schwartz vector whose Fourier transform
has compact support in \((-R,R)\).  Its theta image is zero
for \(x<(qR)^{-1}=\Lambda^{-1}\), by (TB.6).
Such vectors are dense in the parity subspace of
\(\operatorname{Ran}Q\).  Boundedness of \(T\) proves the
operator identity on the whole range.

There is a particularly useful estimate on the complementary
raw time interval.  For \(u>\Lambda\), the integral kernel
of \(T(1-P)\), with input measure \(du\) and output measure
\(dy/y\), is
\[
 K_{\Lambda,g}(y,u)=u^{-1/2}
       \sum_{n>\Lambda u}\chi(n)n^{-1/2}g(ny/u).
 \tag{TB.44}
\]
This follows by writing out the dilation convolution and
changing variables \(u=ny/a\) in its \(n\)-th term.
For compactly supported input the computation is justified
directly; for the additional kernel \(k\), its decay at
infinity makes the resulting \(n\)-sum absolutely convergent.

Put \(F_g(v)=v^{-1/2}g(v)\) and
\[
 A_g(s)=|F_g(s)|+\int_s^\infty |F_g'(v)|\,dv ,
 \tag{TB.45}
\]
where a piecewise derivative suffices for \(k\).
For every \(g\in\mathscr G\), \(A_g\in L^2((0,\infty),ds)\).
For compact smooth \(g\) it is bounded near zero and vanishes
above a fixed bound.  For \(k\), direct calculation gives
\(A_k(s)=1\) for \(0<s\leq1\) and \(A_k(s)=s^{-3/2}\)
for \(s\geq1\).
Abel summation and bounded partial sums of \(\chi\) give
\[
 |K_{\Lambda,g}(y,u)|
       \leq C_\chi\frac{\sqrt y}{u}\,A_g(\Lambda y).
 \tag{TB.46}
\]
Indeed \(n^{-1/2}g(ny/u)=(y/u)^{1/2}F_g(ny/u)\);
the endpoint and total variation above \(n=\Lambda u\)
are exactly bounded by (TB.45) at \(s=\Lambda y\).
It follows that
\[
 \begin{split}
 \|T(1-P)\|_2^2
 &\leq C_\chi
       \int_\Lambda^\infty\frac{du}{u^2}
       \int_0^\infty A_g(\Lambda y)^2\,dy\\
 &\leq C_\chi\Lambda^{-2}\|A_g\|_{L^2(ds)}^2.
 \end{split}
 \tag{TB.47}
\]
The Hilbert--Schmidt kernel extension agrees with the
bounded operator on a dense class of compact raw inputs,
so this is an operator estimate.

As in AC-T27, the projection \(\Pi\) onto the selected raw
prolate vectors satisfies \(\Pi\leq\nu^{-1}PQP\).
Using \(TQ=0\) and then (TB.47),
\[
 \begin{split}
 \sum_j\|T\psi_j\|_2^2
   &=\operatorname{Tr}(T\Pi T^*)\\
   &\leq\nu^{-1}\|TPQ\|_2^2
    =\nu^{-1}\|T(1-P)Q\|_2^2
    \leq C_{\chi,g}(\nu\Lambda^2)^{-1}.
 \end{split}
 \tag{TB.48}
\]
Since each original theta image vanishes above \(\Lambda\),
\(T\psi_j=V(g)(1-S_\Lambda)E_\chi\psi_j\).
This proves (TB.41).
Conjugation of its positive Gram matrix by \(G^{-1/2}\)
and (TB.31) prove (TB.42).

For the last assertion use the classical prolate transition
asymptotic at its center.  In H.J.~Landau,
\emph{An Overview of Time and Frequency Limiting} (1985),
in J.F.~Price (ed.), \emph{Fourier Techniques and Applications},
pp.~201--220, printed pp.~209 and 211,
\url{https://web.eng.ucsd.edu/~massimo/ECE287C/Handouts_files/LandauOverview1985.pdf},
the normalization with angular
bandwidth \(\pi\) has
\[
 \nu_{\lfloor2T+(b/\pi)\log T\rfloor}
                  \longrightarrow(1+e^{\pi b})^{-1}.
 \tag{TB.49}
\]
This is the published scalar asymptotic, used here as
background; its full original proof is not claimed to have
been independently rechecked.  Scaling gives \(T=2c\).
At \(b=0\), therefore, \(\nu_{\lfloor4c\rfloor}\to1/2\).
The largest index of either retained parity is at most
\(\lfloor4c\rfloor\), so its minimum concentration eigenvalue
is at least \(1/4\) for all sufficiently large \(\Lambda\).
Equation (TB.42) now gives (TB.43).
\end{proof}

\begin{theorem}[AC-T30: a positive contraction closes the tested tail comparison]
\label{AC-T30}
Assume \(A\), and retain the actual full parity-index
prolate cutoff.  Define the finite-rank positive contraction
\[
       \widetilde Q_\Lambda=S_\Lambda Q_\Lambda S_\Lambda
              =Z_0Z_0^* .
\]
Then \(0\leq\widetilde Q_\Lambda\leq S_\Lambda\), and for
every \(h\in C_c^\infty(\mathbb R_{>0})\),
\[
 \operatorname{Tr}\bigl((Q_\Lambda-\widetilde Q_\Lambda)V(h)\bigr)
                         \longrightarrow0.
 \tag{TB.50}
\]
For a positive convolution test \(h=g*g^*\), \(g\in\mathscr G\),
and any fixed \(0<\kappa<1\), the error is
\[
 O_{\chi,g,\kappa}
       \bigl(\Lambda^{-1+\kappa}\sqrt{\log\Lambda}\bigr).
 \tag{TB.51}
\]
Consequently the actual renormalized cutoff has asymptotic
positivity:
\[
 \liminf_{\Lambda\to\infty}
 \left(2\log\Lambda\,\|g\|_2^2
       -\operatorname{Tr}(Q_\Lambda V(g*g^*))\right)\geq0.
 \tag{TB.52}
\]

Let
\[
 \widetilde k_{\chi,\Lambda}(\tau)
       =\sum_{j=1}^N|\widehat{(S_\Lambda z_j)}_+(\tau)|^2,
 \qquad z_j=EG^{-1/2}e_j^{\rm std},
\]
where the Fourier transform is in logarithmic coordinates.
Then \(0\leq\widetilde k_{\chi,\Lambda}\leq2\log\Lambda\).
For this original cutoff, T-CUT is now equivalent to the
single remaining positive-measure limit
\[
 \frac{(2\log\Lambda-\widetilde k_{\chi,\Lambda}(\tau))\,d\tau}{2\pi}
     \longrightarrow\sum_\gamma m_{\chi,\gamma}\delta_\gamma
 \quad\hbox{tested against every }\widehat h_+,\quad
                    h\in C_c^\infty(\mathbb R_{>0}).
 \tag{TB.53}
\]
This limit, including its high-frequency control, is not
proved here.  In particular (TB.52) has not been upgraded
to the value of the trace limit.
\end{theorem}
\begin{proof}
The contraction assertion follows from \(0\leq Q_\Lambda\leq1\).
With \(C_g,D_g\) as in AC-T26, the exact tested difference is
\[
 \operatorname{Tr}\bigl((Q_\Lambda-\widetilde Q_\Lambda)
                                      V(g*g^*)\bigr)
       =\|C_g\|_2^2+
                    2\operatorname{Re}\operatorname{tr}(C_g^*D_g).
 \tag{TB.54}
\]
There is no inverse-Gram renormalization term: the compressed
operator is the contraction \(Z_0Z_0^*\), not the orthogonal
projection onto its range.
The energy bound from AC-T26 gives
\(\|D_g\|_2\leq\sqrt{2\log\Lambda}\|g\|_2\).
Apply (TB.43) to \(g^*\), and then Hilbert--Schmidt
Cauchy--Schwarz, to obtain (TB.51).
The test space \(\mathscr G\) is stable under involution.
The explicit factorization \(h=k*j^*\) and polarization in
the remark after AC-T26 therefore prove (TB.50) for every
original smooth compact \(h\).

For \(g\in L^1\cap L^2\), the positive interval pairing is
\[
 \operatorname{Tr}\bigl(V(g)^*S_\Lambda V(g)\bigr)
                    =2\log\Lambda\,\|g\|_2^2,
\]
because \(S_\Lambda V(g)\) is Hilbert--Schmidt.
Since \(S_\Lambda-\widetilde Q_\Lambda\) is positive,
\(\operatorname{Tr}(V(g)^*(S_\Lambda-\widetilde Q_\Lambda)V(g))\)
is nonnegative and finite.
Together with (TB.50)--(TB.51) this proves (TB.52).

The vectors \(z_j\) are orthonormal before compression.
Bessel's inequality against
\(1_{[-L,L]}(t)e^{-i\tau t}\) gives the asserted bound on
\(\widetilde k\).  The finite-rank Fourier trace formula gives
\[
 2\log\Lambda\,h(1)
       -\operatorname{Tr}(\widetilde Q_\Lambda V(h))
 =\frac1{2\pi}\int_{\mathbb R}\widehat h_+(\tau)
                    (2\log\Lambda-\widetilde k_{\chi,\Lambda}(\tau))
                                               \,d\tau .
\]
Combining this identity with (TB.50) proves the equivalence
with (TB.53).  No step identifies the limit measure, so the
last qualification is essential.
\end{proof}

\subsection{The zeros give a lower bound for the remaining measure}

\begin{theorem}[AC-T31: the full zero measure is already a lower bound]
\label{AC-T31}
Assume \(A\) and use the full original parity-index cutoff.
Let
\[
 d\mu_\Lambda(\tau)=
       (2\log\Lambda-\widetilde k_{\chi,\Lambda}(\tau))
                                                     \frac{d\tau}{2\pi}.
\]
For every \(g\in C_c^\infty(\mathbb R_{>0})\),
\[
 \begin{split}
 \liminf_{\Lambda\to\infty}
 \left(2\log\Lambda\,\|g\|_2^2-
                  \operatorname{Tr}(Q_\Lambda V(g*g^*))\right)
 &=
 \liminf_{\Lambda\to\infty}
                  \int|\widehat g_+(\tau)|^2\,d\mu_\Lambda(\tau)\\
 &\geq\sum_\gamma m_{\chi,\gamma}|\widehat g_+(\gamma)|^2.
 \end{split}
 \tag{TB.55}
\]
Every locally finite vague subsequential limit of
\(\mu_\Lambda\), if one exists, dominates the full
zero-counting measure as a positive measure.
No locally finite subsequential limit or matching upper
bound is asserted.  Thus a possible additional positive
measure, or lack of the required boundedness, remains the
precise issue in (TB.53).
\end{theorem}
\begin{proof}
We begin with exact tail evaluations, where the zero
condition makes a new bounded functional available.
Fix a zero \(s_0=1/2+i\gamma\) of multiplicity \(m\).
For \(0\leq j<m\), define on compact raw test vectors
\[
 \ell_{\gamma,j,\Lambda}(\psi)
   =\int_0^{\Lambda^{-1}}E_\chi\psi(x)x^{i\gamma}
                         (\log(\Lambda x))^j\,\frac{dx}{x}.
 \tag{TB.56}
\]
Changing variables \(u=nx\) gives the raw kernel
\[
 H_{\gamma,j,\Lambda}(u)=u^{-1/2+i\gamma}
       \sum_{n>\Lambda u}\chi(n)n^{-1/2-i\gamma}
                                  \bigl(\log(\Lambda u/n)\bigr)^j .
 \tag{TB.57}
\]
For \(u<\Lambda^{-1}\) the sum is over all positive integers.
Expansion of the logarithmic polynomial expresses it as a
linear combination of \(L^{(r)}(s_0,\chi)\), \(0\leq r\leq j\).
It is therefore exactly zero.  These differentiated series
converge on \(\operatorname{Re}s>0\) by bounded character
partial sums and Abel summation.
For \(u\geq\Lambda^{-1}\), the same summation by parts gives
\[
 |H_{\gamma,j,\Lambda}(u)|
                       \leq C_{\chi,\gamma,j}\Lambda^{-1/2}u^{-1}.
 \tag{TB.58}
\]
Indeed the derivative of
\(t^{-1/2-i\gamma}(\log(\Lambda u/t))^j\), integrated over
\(t>\Lambda u\), has total variation at most
\(C_{\gamma,j}(\Lambda u)^{-1/2}\).
The endpoint term has the same bound.

The kernel identity can first be checked with the \(x\)
integral cut off below at \(\varepsilon\).
Abel summation bounds the omitted series tail, after
integration against a bounded compact raw \(\psi\), by
\[
 O_{\psi,\chi,\gamma,j}
       \bigl(\sqrt\varepsilon
                         (1+|\log\varepsilon|)^{j+1}\bigr)=o(1).
\]
This justifies (TB.57) without an absolute interchange of the
original infinite sum.  The same argument applies to
Schwartz raw vectors.  Equations (TB.57)--(TB.58) extend
\(\ell_{\gamma,j,\Lambda}\) boundedly to the raw \(L^2\) space,
with
\[
 \|\ell_{\gamma,j,\Lambda}\|^2\leq C_{\chi,\gamma,j},
 \qquad
 \|\ell_{\gamma,j,\Lambda}(1-P)\|^2
                             \leq C_{\chi,\gamma,j}\Lambda^{-2}.
 \tag{TB.59}
\]
For bandlimited Schwartz vectors the theta image is zero
below \(\Lambda^{-1}\), by Poisson summation.
Density and boundedness therefore give
\(\ell_{\gamma,j,\Lambda}Q=0\).
The projection-order argument of (TB.48) now gives
\[
 \sum_i|\ell_{\gamma,j,\Lambda}(\psi_i)|^2
                    \leq C_{\chi,\gamma,j}(\nu\Lambda^2)^{-1}.
\]
Normalize by the actual Gram matrix and use AC-T28 and
\(\nu\geq1/4\) at the original cutoff.  For every
\(0<\kappa<1\),
\[
 \sum_i\left|
     \int_0^{\Lambda^{-1}}z_i(x)x^{i\gamma}
                         (\log(\Lambda x))^j\,\frac{dx}{x}
                         \right|^2
          =O_{\chi,\gamma,j,\kappa}(\Lambda^{-2+2\kappa}).
 \tag{TB.60}
\]

Each full Mellin transform of \(z_i\) has a zero of order
at least \(m\) at \(s_0\), by (TB.1) and its pre-Plancherel
Mellin identity.  All differentiated integrals converge:
the original compact raw vectors have theta images
\(O(\sqrt x)\) at zero and vanish above \(\Lambda\).
Thus their window integrals in (TB.60) are the negatives
of the displayed tail integrals.
In logarithmic coordinates set
\[
 v_{\gamma,j,L}(t)=
    1_{[-L,L]}(t)e^{-i\gamma t}
          \frac{((t+L)/(2L))^j}{\sqrt{2L}},
 \qquad L=\log\Lambda.
 \tag{TB.61}
\]
Equation (TB.60) shows
\[
 \|Q_\Lambda v_{\gamma,j,L}\|_2^2
      =O_{\chi,\gamma,j,\kappa}
           \bigl(\Lambda^{-2+2\kappa}(2L)^{-2j-1}\bigr)=o(1).
 \tag{TB.62}
\]

Choose any finite set of distinct zero ordinates and include
all the vectors (TB.61) up to their respective multiplicities.
Within one ordinate their Gram matrix is exactly the positive
Hilbert matrix
\[
                 \left(\frac1{j+k+1}\right)_{j,k=0}^{m-1}.
\]
Between distinct fixed ordinates the Gram entries tend to
zero, by integration by parts after \(t=-L+2Ls\).
The full finite Gram matrix therefore tends to an invertible
block diagonal matrix.  Its inverse square root produces an
orthonormal basis \(u_{1,L},\ldots,u_{M,L}\) of this finite
jet space, with
\[
                    \sum_{a=1}^M\|Q_\Lambda u_{a,L}\|_2^2=o(1).
 \tag{TB.63}
\]

Write \(D_\Lambda=S_\Lambda-\widetilde Q_\Lambda\).
On the window, \(0\leq D_\Lambda\leq1\), and
\[
 \|(D_\Lambda^{1/2}-1)u_{a,L}\|_2^2
 \leq\langle u_{a,L},(1-D_\Lambda)u_{a,L}\rangle
                    =\|Q_\Lambda u_{a,L}\|_2^2=o(1).
 \tag{TB.64}
\]
For any fixed \(g\in L^1(dx/x)\cap L^2(dx/x)\), the positive
pairing is finite and
\[
 \int|\widehat g_+|^2\,d\mu_\Lambda
       =\|V(g)^*D_\Lambda^{1/2}\|_2^2
       \geq\sum_{a=1}^M
                  \|V(g)^*D_\Lambda^{1/2}u_{a,L}\|_2^2 .
 \tag{TB.65}
\]
Multiplicative translation by any fixed \(a>0\) changes the
slow polynomial factor in (TB.61) by \(o(1)\) in \(L^2\);
its endpoint change also has norm \(o(1)\).
Consequently
\[
 V(g)^*v_{\gamma,j,L}
            -\overline{\widehat g_+(\gamma)}\,v_{\gamma,j,L}
                         \longrightarrow0
 \tag{TB.66}
\]
in \(L^2\).  For \(g\in L^1\) this follows by dominated
convergence from the fixed-translation assertion.
On each limiting Gram block the scalar in (TB.66) is constant.
The block Gram computation and (TB.64)--(TB.66) give
\[
 \liminf_{\Lambda\to\infty}
       \int|\widehat g_+|^2\,d\mu_\Lambda
       \geq\sum_{\gamma\ {\rm in\ the\ finite\ set}}
                         m_{\chi,\gamma}|\widehat g_+(\gamma)|^2.
 \tag{TB.67}
\]
For compact smooth \(g\), enlarge the finite set to all zeros.
The sum is absolutely convergent by the usual zero count
and rapid decay of \(\widehat g_+\).
Equation (TB.50) proves (TB.55).

Finally (TB.67) also holds when \(\widehat g_+\) is smooth
and compactly supported in frequency, because its inverse
Fourier transform belongs to \(L^1\cap L^2\).
Along any locally finite vague subsequential limit, these
compactly supported squared tests therefore give
\(\mu\geq\sum_\gamma m_{\chi,\gamma}\delta_\gamma\).
For example one may shrink a smooth test equal to one at
each of any finite set of zero ordinates; the resulting
atom bounds, and positivity off that set, give this measure
inequality.  The argument supplies no upper bound on
\(\mu_\Lambda\), as stated.
\end{proof}

\begin{proposition}[AC-T32: the actual adjoint is a co-Poisson map]
\label{AC-T32}
Regard \(E_\chi\) as the closed operator from
\(\mathcal H=L^2(\mathbb R_{>0},du)\) to
\(\mathcal K=L^2(\mathbb R_{>0},dx/x)\) specified in AC-T25.
For \(F\in C_c^\infty(\mathbb R_{>0})\), one has
\[
 (E_\chi^*F)(u)
   =u^{-1/2}\sum_{n\geq1}\overline{\chi(n)}n^{-1/2}F(u/n)
   =\sum_{n\geq1}\frac{\overline{\chi(n)}}n\,g(u/n),
 \qquad g(t)=t^{-1/2}F(t).
 \tag{TB.68}
\]
The parity-\(\epsilon\) extension of this function is Schwartz.
If \(\operatorname{supp}F\subset[a,b]\), it vanishes for
\(0<u<a\), and its raw Fourier transform vanishes for
\(0<|\xi|<(qb)^{-1}\).

Put \(\lambda=(\sqrt q\,\Lambda)^{-1}\), and let
\(\mathscr D_r f(u)=r^{-1/2}f(u/r)\).
Then the closure of
\[
 \left\{\mathscr D_{q^{-1/2}}E_\chi^*F:
                   F\in C_c^\infty((\Lambda^{-1},\Lambda))\right\}
 \tag{TB.69}
\]
is exactly the co-Poisson space generated by
\[
 \sum_{n\geq1}\frac{\overline{\chi(n)}}n\,g(u/n),
       \qquad
 g\in C_c^\infty((\lambda,(q\lambda)^{-1})).
 \tag{TB.70}
\]
Every vector in this closed space and its raw Fourier
transform vanish on \((-\lambda,\lambda)\).
For even characters this is Burnol's
\(W_\lambda^{\overline\chi}\), in his right-Mellin convention;
for odd characters (TB.70) gives the corresponding sine
version directly.
These statements are unconditional.
\end{proposition}
\begin{proof}
For compact raw test vectors the pairing with \(F\) is a
finite sum, and the substitution \(u=nx\) gives
\[
 \langle E_\chi\psi,F\rangle_{\mathcal K}
 =\int_0^\infty\overline{\psi(u)}
       u^{-1/2}\sum_{n\geq1}
          \overline{\chi(n)}n^{-1/2}F(u/n)\,du.
\]
The right side of (TB.68) is locally a finite sum.
It vanishes for \(u<a\).
To check its decay at infinity, put
\(\phi(v)=g(1/v)/v\) for \(v>0\) and extend it with parity
\(\epsilon\).
This is a smooth compactly supported function away from zero,
and (TB.68) equals
\[
             \frac1u\sum_{n\geq1}\overline{\chi(n)}\phi(n/u).
\]
Periodic-character Poisson summation expresses this as a
constant times a sum of \(\widehat\phi(ku/q)\), with the
appropriate finite Fourier coefficients of \(\overline\chi\).
The coefficient at \(k=0\) is zero.
Rapid decay of \(\widehat\phi\), including all its derivatives,
therefore gives rapid decay of the displayed expression and
all its \(u\)-derivatives as \(u\to\infty\).
Together with its vanishing near zero this proves the
Schwartz assertion.  Alternatively the adjoint domain
condition follows immediately from polynomial growth of
\(L(1/2+i\tau,\chi)\) and the rapid decay of the log Fourier
transform of \(F\).
The pairing thus proves (TB.68) for the actual closed adjoint.

Taking adjoints in (TB.6), and then multiplying by the raw
unitary Fourier transform \(\mathcal F_\epsilon\), gives
\[
 \mathcal F_\epsilon E_\chi^*
        =\overline{\omega_\chi}\,E_{\overline\chi}^*J_q,
 \qquad (J_qF)(x)=F((qx)^{-1}).
 \tag{TB.71}
\]
All terms are Schwartz on the present core, so there is no
domain extension in this identity.
The support of \(J_qF\) is contained in
\([(qb)^{-1},(qa)^{-1}]\).
The already proved position-gap statement applied to its
adjoint image proves the Fourier gap.

A direct calculation in (TB.68) gives the exact intertwining
\[
                 \mathscr D_rE_\chi^*=E_\chi^*V_r.
 \tag{TB.72}
\]
For \(r=q^{-1/2}\), \(V_rF\) ranges over all smooth compact
functions in
\[
        (r/\Lambda,r\Lambda)
          =(\lambda,(q\lambda)^{-1}).
 \]
Multiplication by \(t^{-1/2}\) preserves this test-function
space bijectively, so (TB.69) and (TB.70) have exactly the
same generating set.
Position gaps scale by \(r\), while Fourier gaps scale by
\(r^{-1}\); both become \(\lambda\).
Vanishing on a fixed interval is a closed \(L^2\) condition,
so both gaps persist under closure.

For comparison with the source, the left Mellin transform
of (TB.68) is
\[
 \int_0^\infty(E_\chi^*F)(u)u^{s-1}\,du
       =L(1-s,\overline\chi)
                      \int_0^\infty g(u)u^{s-1}\,du.
 \tag{TB.73}
\]
Initially this follows absolutely for \(\operatorname{Re}s<0\);
the Schwartz and gap properties continue it to all \(s\).
Thus the right Mellin convention used in
Burnol, \emph{On Fourier and Zeta(s)}, arXiv:math/0112254v6,
Section 6, the definitions preceding Theorems 6.29--6.30
(printed pp.\ 44--46), gives multiplication by
\(L(s,\overline\chi)\), exactly as asserted.
Allowing smooth generators flat at the interval endpoints
gives the same closure: cut them off a positive distance
from the endpoints and use convergence in every smooth
norm, hence in the adjoint graph norm.
\end{proof}

\begin{theorem}[AC-T33: a uniform raw co-Poisson--prolate bound]
\label{AC-T33}
Let \(\mathscr G=C_c^\infty(\mathbb R_{>0})+\mathbb Ck\)
be the test class of AC-T29, where
\(k(e^t)=\tfrac12e^{-|t|}\).
For \(0<a<b\), put \(S_{a,b}=1_{[a,b]}\).
For every \(g\in\mathscr G\),
\[
 \bigl\|1_{(b,\infty)}E_\chi^*S_{a,b}V(g)\bigr\|_2
                   \leq C_{\chi,g},
 \tag{TB.74}
\]
where the Hilbert--Schmidt constant is independent of both
\(a\) and \(b\).
The full operator is Hilbert--Schmidt, with
\[
       \|E_\chi^*S_{a,b}V(g)\|_2^2
                   \leq C_{\chi,g}^2(1+\log(b/a)).
 \tag{TB.75}
\]

Let \(\Pi_\Lambda\) be the orthogonal projection in the
unweighted raw space onto the actual parity prolate family
with indices \(n\leq4\Lambda^2/q\).
Then, for all sufficiently large \(\Lambda\),
\[
 \bigl\|(1-\Pi_\Lambda)
               E_\chi^*S_\Lambda V(g)\bigr\|_2
                     \leq C'_{\chi,g}.
 \tag{TB.76}
\]
In particular this is a uniform Hilbert--Schmidt bound
for the discarded part of the actual co-Poisson adjoint,
with the original cutoff threshold.
It does not assume A.
\end{theorem}
\begin{proof}
Set
\[
 F_g(v)=v^{-1/2}g(v^{-1}),\qquad
 B_g(s)=|F_g(s)|+\int_s^\infty|F_g'(v)|\,dv.
 \tag{TB.77}
\]
The function \(F_g\) is locally absolutely continuous and
\(B_g\in L^2(0,\infty)\).
For compact smooth \(g\) this follows from compact support
away from zero.
For \(g=k\), direct integration gives
\[
                B_k(s)=
                 \begin{cases}1,&0<s\leq1,\\
                               s^{-3/2},&s\geq1.
                  \end{cases}
\]
The assertion for their finite linear combinations follows
from the triangle inequality.

By (TB.68), the kernel of \(E_\chi^*S_{a,b}V(g)\),
with measures \(du\) and \(dy/y\), is
\[
 K_{a,b,g}(u,y)
   =u^{-1/2}\!
       \sum_{u/b\leq n\leq u/a}
          \overline{\chi(n)}n^{-1/2}g(u/(ny)).
 \tag{TB.78}
\]
This sum is finite for every \(u>0\).
The formula first holds on compact smooth inputs.
Such inputs after \(S_{a,b}V(g)\) have piecewise smooth
logarithmic profiles, so their Fourier transforms decay
as \(O((1+|\tau|)^{-1})\).
The polynomial bound of exponent \(5/16<1/2\) in AC-T25
places them in the domain of \(E_\chi^*\).
Thus (TB.78) is the kernel of the closed-adjoint composition
on this dense core.

For \(u>b\), Abel summation using bounded character partial
sums gives
\[
 |K_{a,b,g}(u,y)|
    \leq C_\chi\,\frac{\sqrt y}{u}\,B_g(y/b).
 \tag{TB.79}
\]
Indeed
\(n^{-1/2}g(u/(ny))=(y/u)^{1/2}F_g(ny/u)\).
The first retained argument is at least \(y/b\);
both endpoint terms and the total variation on the
retained interval are bounded by a fixed multiple of
\(B_g(y/b)\).
Consequently
\[
 \int_b^\infty\!\!\int_0^\infty
       |K_{a,b,g}(u,y)|^2\,\frac{dy}{y}\,du
 \leq C_\chi^2
       \left(\int_b^\infty\frac{du}{u^2}\right)
       \left(\int_0^\infty B_g(y/b)^2\,dy\right)
 =C_\chi^2\|B_g\|_2^2.
 \tag{TB.80}
\]
The kernel is zero for \(u<a\).
When \(a\leq u\leq b\), the first retained integer is \(1\).
The same Abel bound gives
\[
 |K_{a,b,g}(u,y)|
       \leq C_\chi\,\frac{\sqrt y}{u}\,B_g(y/u).
 \]
Its squared integral in \(y\) is at most
\(C_\chi^2u^{-1}\|B_g\|_2^2\).
Integration over \(a\leq u\leq b\), together with
(TB.80), proves (TB.74)--(TB.75), and extends the
composition canonically to a Hilbert--Schmidt operator.

Put \(P=1_{[-\Lambda,\Lambda]}\) in the parity raw space
and \(Q=\mathcal F_\epsilon^*1_{[-\Lambda/q,\Lambda/q]}
                                      \mathcal F_\epsilon\).
Equation (TB.74) bounds
\((1-P)E_\chi^*S_\Lambda V(g)\) uniformly.
For the Fourier tail use (TB.71) and
\[
 J_qS_\Lambda=S_{1/(q\Lambda),\,\Lambda/q}J_q,\qquad
 J_qV(g)=V(g^\iota)J_q,\quad g^\iota(x)=g(x^{-1}).
 \tag{TB.81}
\]
The class \(\mathscr G\) is stable under \(g\mapsto g^\iota\).
Applying (TB.74) to \(\overline\chi\), upper endpoint
\(b=\Lambda/q\), and the interval in (TB.81) gives a
uniform bound for
\((1-Q)E_\chi^*S_\Lambda V(g)\).

At the original full cutoff, the largest discarded
eigenvalue \(\eta_\Lambda\) of \(PQP\) is at most \(3/4\)
for all sufficiently large \(\Lambda\), by the same
Landau transition asymptotic used in AC-T29.
For every raw \(f\), the spectral theorem on
\(\operatorname{ran}P\) gives
\[
 \|(P-\Pi_\Lambda)f\|^2
   \leq(1-\eta_\Lambda)^{-1}\|(1-Q)Pf\|^2.
 \]
Since \(\Pi_\Lambda\leq P\), it follows that
\[
 \|(1-\Pi_\Lambda)f\|^2
      \leq9\|(1-P)f\|^2+8\|(1-Q)f\|^2.
 \tag{TB.82}
\]
Apply this inequality to the columns of
\(E_\chi^*S_\Lambda V(g)\) in any orthonormal basis of
the target space and sum.
The two uniform tail estimates prove (TB.76).
\end{proof}

\begin{remark}[The remaining conversion to the polar defect]
The bound (TB.76) concerns the raw adjoint before taking its
polar normalization.
The positive measure (TB.53) instead involves
\(D_\Lambda=S_\Lambda-S_\Lambda Q_\Lambda S_\Lambda\),
where \(Q_\Lambda=E\,(E^*E)^{-1}E^*\) for the selected
finite family.
No bounded inverse for multiplication by
\(L(1/2+i\tau,\chi)\) on the whole critical line is available:
that multiplier has the stipulated critical zeros.
Thus (TB.76) is an actual uniform estimate for the relevant
co-Poisson family, but its conversion to a uniform tested
trace bound for \(D_\Lambda\), and then to the matching
upper limit in (TB.53), still requires proof.
Burnol's Theorem 6.30 is a completeness statement in the
Sonine spaces; it does not itself perform this conversion
for \(\Pi_\Lambda\).
\end{remark}

\begin{theorem}[AC-T34: compact boundary limits and a raw nuclear bound]
\label{AC-T34}
For \(0<r<1\), define
\[
 B_{\chi,g}^{(r)}
       =1_{(1,\infty)}E_\chi^*1_{[r,1]}V(g):
                       \mathcal K\longrightarrow\mathcal H .
 \tag{TB.83}
\]
For \(g\in\mathscr G\) these operators converge in
Hilbert--Schmidt norm to the operator with kernel
\[
 K_{\chi,g}^{(0)}(u,y)
     =1_{u>1}\,u^{-1/2}
          \sum_{n\geq u}\overline{\chi(n)}n^{-1/2}g(u/(ny)),
 \tag{TB.84}
\]
and
\[
       \|B_{\chi,g}^{(r)}-B_{\chi,g}^{(0)}\|_2
                         \leq C_{\chi,g}\sqrt r .
 \tag{TB.85}
\]
Here the raw output measure is \(du\) and the input
measure is \(dy/y\).
For any \(0<a<b\), the actual raw upper boundary operator
of AC-T33 satisfies the exact unitary scaling identity
\[
 \mathscr D_b^*
       \bigl(1_{(b,\infty)}E_\chi^*S_{a,b}V(g)\bigr)V_b
                         =B_{\chi,g}^{(a/b)}.
 \tag{TB.86}
\]

If \(g\in C_c^\infty(\mathbb R_{>0})\), then for every
positive integer \(m\) there is a constant \(C_{\chi,g,m}\)
such that, for all \(j\geq1\) and \(0<r<1\),
\[
 \begin{split}
 s_j(B_{\chi,g}^{(0)})&\leq C_{\chi,g,m}j^{-m},\\
 s_j(B_{\chi,g}^{(r)}-B_{\chi,g}^{(0)})
                          &\leq C_{\chi,g,m}\sqrt r\,j^{-m}.
 \end{split}
 \tag{TB.87}
\]
In particular the convergence in (TB.85) also holds in
trace norm, at the same rate.
For the actual full prolate family this yields
\[
 s_j\bigl((1-\Pi_\Lambda)
                 E_\chi^*S_\Lambda V(g)\bigr)
                     \leq C'_{\chi,g,m}j^{-m}
 \tag{TB.88}
\]
for every \(j\geq1\), uniformly for all sufficiently large
\(\Lambda\).
Thus this raw discarded operator is uniformly trace class,
and uniformly belongs to every Schatten class of positive
index.  The theorem is unconditional.
\end{theorem}
\begin{proof}
The sum in (TB.84) converges for each \(u,y>0\).
For compact \(g\) it has only finitely many nonzero terms;
for the Green kernel its summands are \(O_{u,y}(n^{-3/2})\).
The Abel estimate (TB.79) with upper raw endpoint \(1\)
bounds its absolute value by
\[
       C_\chi\,1_{u>1}\frac{\sqrt y}{u}B_g(y).
 \tag{TB.89}
\]
This is a square-integrable majorant for the kernel
measures just specified.
The difference between (TB.84) and the kernel for
\(B_{\chi,g}^{(r)}\) is the sum over \(n>u/r\).
Abel summation at this new lower endpoint gives
\[
 |K_{\chi,g}^{(r)}(u,y)-K_{\chi,g}^{(0)}(u,y)|
       \leq C_\chi\,1_{u>1}\frac{\sqrt y}{u}B_g(y/r).
 \tag{TB.90}
\]
Integrating its square gives
\[
 \|B_{\chi,g}^{(r)}-B_{\chi,g}^{(0)}\|_2^2
 \leq C_\chi^2
       \int_1^\infty\frac{du}{u^2}
       \int_0^\infty B_g(y/r)^2\,dy
 =C_\chi^2r\|B_g\|_2^2 .
 \]
This proves (TB.85).
Equation (TB.72), the equality
\(S_{a,b}V_b=V_bS_{a/b,1}\), and commutativity of
multiplicative convolutions with \(V_b\) prove
(TB.86), including all normalization factors.

We now assume \(g\) compact smooth.
Identify the input space \(\mathcal K\) with \(L^2(ds)\)
by \(y=e^s\), and let
\[
                         H=1-\partial_s^2+s^2 .
 \]
Its eigenvalues are \(2j+2\), \(j\geq0\).
For every \(j\geq0\), put
\(g_j(x)=(x\,d/dx)^jg(x)\).
Differentiating a kernel \(j\) times in \(s\) simply
replaces \(g\) by \((-1)^jg_j\).
All \(g_j\) are compact smooth.
Each \(B_{g_j}(e^s)\) is bounded as \(s\to-\infty\)
and zero for sufficiently large \(s\).
Consequently, for all nonnegative integers \(k,j\),
\[
 \int_{\mathbb R}
      |s|^{2k}e^s B_{g_j}(e^s)^2\,ds<\infty.
 \tag{TB.91}
\]
The derivative and weighted versions of (TB.89)
therefore show that applying any power of \(H\) in the
input variable leaves a square-integrable kernel.
Integration by parts on compact input tests gives
\[
             B_{\chi,g}^{(0)}H^m\in\mathcal S_2,\qquad
             \|B_{\chi,g}^{(0)}H^m\|_2\leq C_{\chi,g,m}.
 \tag{TB.92}
\]
This refers to the bounded extension of the composition
initially defined on the oscillator core.

For the difference put
\(C_r=B_{\chi,g}^{(r)}-B_{\chi,g}^{(0)}\).
Right composition by the unitary \(V_r\) changes its
kernel to \(K_{\chi,g}^{(r)}(u,ry)-K_{\chi,g}^{(0)}(u,ry)\).
Equation (TB.90) then gives
\[
 |(C_rV_r)(u,e^s)|
       \leq C_\chi\sqrt r\,1_{u>1}
                           \frac{e^{s/2}}u B_g(e^s).
 \tag{TB.93}
\]
Differentiation in \(s\) again replaces \(g\) by a log
derivative, with no derivative of a summation endpoint:
the condition \(n>u/r\) is independent of \(s\).
The weighted integrals (TB.91) now give
\[
                      \|(C_rV_r)H^m\|_2
                            \leq C_{\chi,g,m}\sqrt r.
 \tag{TB.94}
\]
Factor each operator in (TB.92)--(TB.94) as
\((BH^m)H^{-m}\).
Since \(BH^m\) is bounded and
\(s_j(H^{-m})=(2j)^{-m}\), the elementary product
singular-value inequality gives (TB.87).
Taking \(m=2\) and summing over \(j\) proves the trace-norm
rate.  Higher \(m\) gives every positive Schatten index.

It remains to pass from the two boundaries to the actual
discarded prolate component, without changing its threshold.
Let
\[
 T_\Lambda f=
        \bigl(3(1-P)f,\sqrt8(1-Q)f\bigr)
                      \in\mathcal H\oplus\mathcal H,
 \]
where \(P,Q\) are the raw projections in AC-T33.
Inequality (TB.82) implies the factorization
\[
               1-\Pi_\Lambda=R_\Lambda T_\Lambda,
                    \qquad \|R_\Lambda\|\leq1.
 \tag{TB.95}
\]
For clarity, define
\(R_\Lambda(T_\Lambda f)=(1-\Pi_\Lambda)f\) on the
range of \(T_\Lambda\).
The same inequality proves that this is well defined and
contractive; extend by continuity to its closed range and
by zero on the orthogonal complement.

The first component of
\(T_\Lambda E_\chi^*S_\Lambda V(g)\) is, up to the
displayed constant and unitary dilations, \(B_{\chi,g}^{(r)}\)
with \(r=\Lambda^{-2}\).
By (TB.71) and (TB.81), its second component is, up to
unitaries and a unit scalar,
\(B_{\overline\chi,g^\iota}^{(r)}\), with the same \(r\).
Equation (TB.87), and the singular-value inequality for a
sum, give \(s_j\leq C_mj^{-m}\) for each component,
uniformly in \(\Lambda\).
Approximate each component by a rank \(n\) singular-value
truncation.  Their column operator then has a rank at most
\(2n\) approximation with operator-norm error at most
\(C_m(n+1)^{-m}\).
Composition with the contraction \(R_\Lambda\) does not
increase that error or its rank.
This proves (TB.88), after adjusting the constants for
the indices \(2n+1\).
\end{proof}

\begin{remark}[Order of limits in a prospective upper argument]
AC-T34 provides compact boundary operators with quantitative
trace-norm convergence before arithmetic polar normalization.
It does not identify the positive defect \(D_\Lambda\)
with a Sonine projection.
Under A, one cannot replace the finite zero-jet spaces
used in AC-T31 by the span of all zero jets at a fixed log
window: their closed span is already the whole window
space, by the compact-support zero-count argument of
AC-T13.  A prospective completeness argument must preserve
the order of limits: a fixed finite zero set first,
\(\Lambda\to\infty\) next, and enlargement of the zero set
afterwards.
\end{remark}

\begin{theorem}[AC-T35: subpower control of the truncated-image Gram]
\label{AC-T35}
Assume A and retain the original full prolate family.
Let \(G^0=G-H\) be the Gram matrix of the actual vectors
\(S_\Lambda E_\chi\psi_j\).
For every \(\kappa>0\) there are constants
\(c_{\chi,\kappa}>0\) and \(\Lambda_{\chi,\kappa}\) such that
\[
 G^0\geq c_{\chi,\kappa}\Lambda^{-\kappa}I,
 \qquad
 \|(G^0)^{-1}\|\leq c_{\chi,\kappa}^{-1}\Lambda^\kappa
       \quad(\Lambda\geq\Lambda_{\chi,\kappa}).
 \tag{TB.96}
\]
The assertion concerns the sharply truncated theta images,
not just the full Gram matrix in AC-T28.
\end{theorem}
\begin{proof}
We give the finite Dirichlet-polynomial estimate first,
including its dependence on all vertical frequencies.
For arbitrary \(\alpha,\beta>0\), A implies
\[
 \left|\sum_{n\leq N}
      \frac{\mu(n)\chi(n)}{n^{1/2+i\tau}}\right|
       \leq C_{\chi,\alpha,\beta}
                    (1+N)^\alpha(2+|\tau|)^\beta
       \quad(N\geq1,\ \tau\in\mathbb R).
 \tag{TB.97}
\]
Indeed the zero-free logarithm argument in AC-T28, with
its disks uniformly covering each smaller right half-plane,
gives
\[
 |L(\sigma+it,\chi)^{-1}|
       \leq C_{\chi,\delta,\eta}(2+|t|)^\eta
       \quad(\sigma\geq1/2+\delta)
 \tag{TB.98}
\]
for every fixed \(\delta,\eta>0\).
For \(\sigma\geq2\) the Euler series gives the bound directly.

For \(N\geq2\), set \(X=N+1/2\),
\(c=1/2+1/\log X\), \(s=1/2+i\tau\), and
\(T=X^3(2+|\tau|)^3\).
Truncated Perron inversion gives
\[
 \sum_{n\leq N}\mu(n)\chi(n)n^{-s}
   =\frac1{2\pi i}\int_{c-iT}^{c+iT}
        \frac{X^w}{wL(s+w,\chi)}\,dw
           +O\!\left(\frac{\sqrt X\log X}{T}\right).
 \tag{TB.99}
\]
For completeness the truncation bound follows termwise
from the inverse Mellin kernel of \(1_{(1,\infty)}\);
integration by parts in its omitted vertical tails bounds
the \(n\)-th error by
\[
 C n^{-1/2}(X/n)^c
             \min\{1,(T|\log(X/n)|)^{-1}\}.
 \]
For \(X/2\leq n\leq2X\), use
\(|\log(X/n)|\geq C|X-n|/X\) and
\(|X-n|\geq1/2\); the harmonic sum is \(O(\log X)\).
For \(n<X/2\) and \(n>2X\), use
\(|\log(X/n)|\geq\log2\) and
\(\sum n^{-1-1/\log X}\leq C\log X\).
These three estimates prove the error in (TB.99),
uniformly in \(\tau\).

Choose \(0<\delta<c\).
Move the finite rectangular contour to
\(\operatorname{Re}w=\delta\).
There is no pole between the two lines: \(w=0\) is to
their left and (TB.98) applies to \(1/L(s+w,\chi)\).
The new vertical integral is at most
\[
 C_{\chi,\delta,\eta}
         X^\delta(2+|\tau|+T)^\eta\log(2T).
 \]
Each horizontal segment is at most
\[
 C_{\chi,\delta,\eta}
         X^c(2+|\tau|+T)^\eta/T.
 \]
Since \(X^c=e\sqrt X\) and \(T\geq X^3(2+|\tau|)^3\),
choose \(\delta,\eta\) sufficiently small in terms of
\(\alpha,\beta\), and absorb the logarithm in arbitrarily
small positive powers.
The horizontal terms and the Perron error are bounded,
and the vertical term is bounded by the right side of
(TB.97).  This proves (TB.97) for every real \(\tau\),
not merely on a fixed frequency interval.

We next need an unconditional raw concentration fact.
There is \(d_\chi>0\) such that, for every vector \(\psi\)
in the selected prolate span and all sufficiently large
\(\Lambda\),
\[
           \int_{\Lambda^{-1}}^\Lambda|\psi(u)|^2\,du
                         \geq d_\chi\|\psi\|_{\mathcal H}^2.
 \tag{TB.100}
\]
Here is a proof that retains the fixed conductor.
Write \(R=\Lambda/q\), \(c_0=\Lambda R\), and dilate the
raw variable by \(R\).
The time interval becomes \([-c_0,c_0]\), the Fourier
band becomes \([-1,1]\), and
\([-\Lambda^{-1},\Lambda^{-1}]\) becomes
\([-1/q,1/q]\).
Let \(P_{c_0}\) denote the new time projection, let
\(Q_1\) denote the fixed band projection, and let
\(\Pi_{c_0}\) be the dilated selected projection.
Then \(P_{c_0}Q_1P_{c_0}\to Q_1\) strongly.
Its spectral projections at any fixed level in \((0,1)\)
also converge strongly to \(Q_1\): approximate the step
function above and below by continuous functions that
agree on the spectrum \(\{0,1\}\) of the limit.
The retained and discarded eigenvalue bounds used in
AC-T29 and AC-T33 sandwich \(\Pi_{c_0}\) between such
projections, so \(\Pi_{c_0}\to Q_1\) strongly.

Put \(A_0=1_{[-1/q,1/q]}\) in the chosen parity space.
For all large \(c_0\),
\[
 0\leq A_0\Pi_{c_0}A_0
      \leq4A_0Q_1A_0.
 \tag{TB.101}
\]
The operator \(K_q=A_0Q_1A_0\) is compact, indeed trace
class, since the band projection has a continuous kernel
on the bounded interval.
Its norm \(\kappa_q\) is strictly less than one.
Otherwise compactness would produce a unit vector both
time supported in that interval and bandlimited, contrary
to the elementary analytic-continuation uncertainty
argument in AC-T10.
Strong convergence in (TB.101), together with compact
domination, gives norm convergence to \(K_q\).
To see this directly, compress by a finite-rank spectral
projection of \(K_q\).  On its complement the positive
operators in (TB.101) have uniformly small norm; their
off-diagonal blocks are small by Cauchy--Schwarz for
positive operators, while convergence on the finite
block is uniform.
It follows that
\(\|A_0\Pi_{c_0}A_0\|\leq(1+\kappa_q)/2<1\)
eventually.  This proves (TB.100), with
\(d_\chi=(1-\kappa_q)/2\), also after restriction to the
positive half-line.

Fix once and for all \(0<s_0<1/2\).
Set \(L=\log\Lambda\), and for \(t<L\) put
\[
             v(t)=e^{t/2}\psi(e^t).
 \]
Then
\[
 \|v\|_2=\|\psi\|_{\mathcal H},\qquad
 \|v'\|_2\leq C_\chi(1+\Lambda^2)\|\psi\|_{\mathcal H}.
 \tag{TB.102}
\]
Both norms in (TB.102) are taken on \((-\infty,L)\);
the endpoint cuts are treated in the next estimate.
This follows from the prolate derivative bound used in
AC-T28, because the retained concentration values are
bounded below.
For any subinterval \(J\subset(-\infty,L)\), zero extension
of \(v1_J\) satisfies
\[
 \|v1_J\|_{H^{s_0}(\mathbb R)}
      \leq C_{s_0}
        \|v\|_2^{1-s_0}(\|v\|_2+\|v'\|_2)^{s_0}.
 \tag{TB.103}
\]
The constant is independent of both endpoints.
This is the two-endpoint version of the Fourier proof
in AC-T28: integration by parts gives the transform of
\(v'1_J\) divided by \(i\tau\), and at most two endpoint
terms divided by \(i\tau\).
At every endpoint,
\(|v(t)|^2\leq2\|v\|_2\|v'\|_2\), by integrating from
\(-\infty\).
Split the Fourier integral at
\(1+\|v'\|_2/\|v\|_2\).
The endpoint-tail integral converges exactly because
\(s_0<1/2\), and gives (TB.103).
Thus discontinuous interval cutoffs are included.

Let
\[
 e_0(t)=1_{[-L,L]}(t)(E_\chi\psi)(e^t),
               \qquad N=\lfloor\Lambda^2\rfloor .
 \]
On this log window the theta sum is finite:
\[
 e_0(t)=\sum_{n\leq N}\chi(n)n^{-1/2}
        1_{[-L,L-\log n]}(t)v(t+\log n).
 \tag{TB.104}
\]
Each summand before its coefficient is a translate of
a zero-extended interval restriction of \(v\).
Equations (TB.102)--(TB.103), and
\(\sum_{n\leq N}n^{-1/2}\leq2\sqrt N\), imply the
deliberately coarse polynomial bound
\[
 \|e_0\|_{H^{s_0}(\mathbb R)}
          \leq C_{\chi,s_0}
                 \Lambda^{1+2s_0}\|\psi\|_{\mathcal H}.
 \tag{TB.105}
\]

Finite Möbius inversion, with all functions extended by
zero outside the window, gives
\[
 1_{[-L,L]}(t)v(t)=
 1_{[-L,L]}(t)\sum_{n\leq N}\mu(n)\chi(n)n^{-1/2}
                                      e_0(t+\log n).
 \tag{TB.106}
\]
This is exact almost everywhere: substituting the theta
sum groups each integer \(k=mn\) with coefficient
\(\chi(k)\sum_{n\mid k}\mu(n)\), which is zero unless
\(k=1\).  Vanishing of \(\chi\) at ramified integers is
compatible with this complete multiplicativity.
The full finite convolution in (TB.106) has Fourier
multiplier equal to the polynomial in (TB.97).
For any \(0<\beta<s_0\), set \(d=\beta/s_0\).
Plancherel, (TB.97), Sobolev interpolation, (TB.100),
and (TB.105) give
\[
 \begin{split}
 \sqrt{d_\chi}\,\|\psi\|_{\mathcal H}
 &\leq C_{\chi,\alpha,\beta}N^\alpha
                                         \|e_0\|_{H^\beta}\\
 &\leq C_{\chi,\alpha,\beta,s_0}
      \Lambda^{2\alpha+d(1+2s_0)}
               \|e_0\|_2^{\,1-d}\|\psi\|_{\mathcal H}^{\,d}.
 \end{split}
 \tag{TB.107}
\]
Consequently
\[
 \|e_0\|_2^2
    \geq C_{\chi,\alpha,\beta,s_0}
       \Lambda^{-\{4\alpha+2d(1+2s_0)\}/(1-d)}
                                      \|\psi\|_{\mathcal H}^2.
 \tag{TB.108}
\]
Choose \(\alpha,d>0\) so small that the displayed exponent
is at most the prescribed \(\kappa\).
The raw prolate basis is orthonormal, so (TB.108) is
exactly the matrix inequality (TB.96).
\end{proof}

\begin{proposition}[AC-T36: exact asymmetric dilation and its metric constants]
Let \(D_r\psi(u)=r^{-1/2}\psi(u/r)\) on the raw parity
space, and let \(V_rF(x)=F(x/r)\) on
\(L^2(\mathbb R_+,d^*x)\).
For \(r>0\), put
\[
 X=r\Lambda,\qquad Y=\Lambda/r,\qquad
 \mathcal I_{X,Y}=[1/Y,X].
 \tag{TB.109}
\]
Dilation carries the actual raw time and frequency cutoffs
\((\Lambda,\Lambda/q)\) to \((X,Y/q)\), keeps their product
\(\Lambda^2/q\) and the selected prolate indices unchanged,
and carries the actual theta-image projection to
\[
 Q_{X,Y}=V_rQ_\Lambda V_r^*,\qquad
 S_{X,Y}=V_rS_\Lambda V_r^*
          =1_{\mathcal I_{X,Y}}.
 \tag{TB.110}
\]
Consequently, for every compact smooth norm test \(h\),
\[
 \operatorname{Tr}(Q_{X,Y}V(h))
       =\operatorname{Tr}(Q_\Lambda V(h)),\qquad
 \operatorname{Tr}(S_{X,Y}Q_{X,Y}S_{X,Y}V(h))
       =\operatorname{Tr}(S_\Lambda Q_\Lambda S_\Lambda V(h)).
 \tag{TB.111}
\]
These are exact equalities, including the full arithmetic
Gram normalization.

Suppose \(X,Y\geq1\) and \(\alpha\geq0\), and equip the target log line with
weight \(w_\alpha(t)=e^{2\alpha|t|}\).
For any finite column map \(F\) whose columns are supported
in \(\mathcal I_{X,Y}\), set
\(G_0=F^*F\) and \(G_\alpha=F^*w_\alpha F\), where the
adjoints in this display use the unweighted norm.
Writing \(M=\max(X,Y)\), one has the exact comparisons
\[
 G_0\leq G_\alpha\leq M^{2\alpha}G_0,\qquad
 M^{-2\alpha}G_0^{-1}\leq G_\alpha^{-1}\leq G_0^{-1}
 \tag{TB.112}
\]
when the columns are independent.
The identity map between the unweighted and weighted
window spaces has norm at most \(M^\alpha\), and its
inverse has norm at most one.

In particular, the choice \(r=\Lambda^\theta\),
\(0<\theta<1\), gives \(X,Y\to\infty\) and \(Y/X\to0\)
without changing (TB.111).
If \(\alpha=\alpha_\Lambda\to0\), all the factors
\(M^{O(\alpha_\Lambda)}\) are \(\Lambda^{o(1)}\).
For example, \(\alpha_\Lambda=(\log\Lambda)^{-1/2}\)
also satisfies \(\alpha_\Lambda\log X\to\infty\).
\end{proposition}
\begin{proof}
The Fourier scaling identity sends the band
\([-\Lambda/q,\Lambda/q]\) to \([-Y/q,Y/q]\), and
the position support becomes \([-X,X]\).
It conjugates the product of the two raw projections
unitarily, so its eigenvalues and the chosen index
threshold are unchanged.
Direct substitution in the theta sum gives
\[
 E_\chi D_r=V_rE_\chi.
 \tag{TB.113}
\]
The identity extends from the smooth core to all the
finite raw spaces in question by the closed Mellin
realization of \(E_\chi\).
Thus the Gram matrices of the dilated columns equal
the original matrices, proving (TB.110).
Multiplicative convolution commutes with \(V_r\);
unitary trace invariance proves (TB.111).
On the stated window,
\(1\leq e^{2\alpha|\log x|}\leq M^{2\alpha}\).
Integration gives the first pair of positive-matrix
inequalities in (TB.112); inversion reverses order and
gives the second pair.
The remaining assertions follow by substitution.
\end{proof}

\begin{theorem}[AC-T37: a shrinking-strip causal slow-band estimate under \(\mathsf A\)]
Fix a nonprincipal primitive Dirichlet character, assume
\(\mathsf A\), and list its nontrivial zeros with their
full multiplicities as \(1/2+i\gamma_j\).
Fix \(0<\alpha_0<1/2\).
For \(0<\alpha\leq\alpha_0\), let \(B_\alpha\) be the
pure Hardy Blaschke product with zeros
\(z_{j,\alpha}=\alpha-i\gamma_j\), let \(P_\alpha\)
be the model-space projection on the time half-line,
and let \(V_\alpha\) be its inner isometry.
These are the actual objects of MA-CGP.12 and MA-SB.1.

Let \(g\in C_c^\infty((0,\infty))\), with support in
\((0,b)\), and put \(k_\alpha(t)=e^{\alpha t}g(t)\).
Define the actual trace-norm tails
\[
 \omega_{\alpha,g}(R)
   =\|1_{[R,\infty)}C_{k_\alpha}^+P_\alpha\|_1,\qquad
 \nu_{\alpha,g}(R)
   =\|1_{[R,\infty)}H_{B_\alpha k_\alpha}\|_1.
 \tag{TB.114}
\]
There is \(C_{\chi,g,\alpha_0}\), independent of
\(\alpha,R,N\), such that for \(R>b\) and \(N\geq1\),
\[
 \omega_{\alpha,g}(R)+\nu_{\alpha,g}(R)
 \leq C_{\chi,g,\alpha_0}
       \left\{\left(\sum_{j>N}\frac1{1+\gamma_j^2}\right)^{1/2}
                        +\mathcal H_N(\alpha(R-b))\right\},
 \tag{TB.115}
\]
where the explicit finite function
\[
 \begin{split}
 p_j(v)&=\sum_{\ell=0}^{j-1}
              \binom{j-1}{\ell}\frac{(2v)^\ell}{\ell!},\\
 \mathcal H_N(a)&=
       \left(2\sum_{j=1}^N
                \int_a^\infty e^{-2v}p_j(v)^2\,dv\right)^{1/2}
       \leq C_N(1+a)^{N-1}e^{-a}
 \end{split}
 \tag{TB.116}
\]
tends to zero as \(a\to\infty\).
No lower bound on distinct zero gaps enters the estimate.

Consequently, the actual weighted PW trace in
MA-SB.15 has its stated zero-sum limit along every
source-thin path with
\[
 X,Y\longrightarrow\infty,\qquad Y/X\longrightarrow0,
 \qquad 0<\alpha_{X,Y}\leq\alpha_0,\qquad
             \alpha_{X,Y}\log X\longrightarrow\infty,
 \tag{TB.117}
\]
for these causal tests, even if \(\alpha_{X,Y}\to0\).
In logarithmic notation its limit is
\[
 \operatorname{Tr}_{H_{\alpha_{X,Y}}}
       (M_{\rho_{X,\delta}}C_gQ^{\,w}_{-\log Y})
       \longrightarrow
       -\sum_j\int_0^\infty g(t)e^{i\gamma_jt}\,dt.
 \tag{TB.118}
\]
Here \(Q^{\,w}_{-\log Y}\) is the completed weighted
PW projection, not the finite projection in (TB.110).
\end{theorem}
\begin{proof}
All Hardy norms use boundary measure \(dt/(2\pi)\).
Let \(e_{j,\alpha}\) be the normalized Malmquist basis,
in the specified order, and let \(P_{N,\alpha}\) be
the projection onto its first \(N\) vectors.
The elementary model-space decomposition used in
MA-GP.89 proves that these vectors form an orthonormal
basis, including repeated zeros.
The modulus of each preceding Blaschke factor is one
on the boundary.  For
\(r(z)=(1+z)^{-1}\), the Poisson calculation in
MA-GP.90--92 therefore gives
\[
 \|M_r e_{j,\alpha}\|_2^2
  =\frac1\pi\int_{\mathbb R}
       \frac{\alpha\,dt}
            {(1+t^2)((t+\gamma_j)^2+\alpha^2)}
  \leq \frac{C_{\alpha_0}}{1+\gamma_j^2}.
 \tag{TB.119}
\]
For clarity, this bound is uniform down to \(\alpha=0\):
if \(|\gamma_j|\geq2\), split the integral into
\(|t+\gamma_j|\leq|\gamma_j|/2\), where
\((1+t^2)^{-1}\leq C/\gamma_j^2\), and its
complement, where the Poisson density is at most
\(C\alpha_0/\gamma_j^2\).
Bounded \(\gamma_j\) follows from total Poisson mass.
The standard zero count ensures
\(\sum_j(1+\gamma_j^2)^{-1}<\infty\).

Write \(T_{r,\alpha}=P_\alpha M_r|_{K_{B_\alpha}}\)
and \(\mathcal R_\alpha=T_{r,\alpha}^*\), extended
by zero on the orthogonal complement.
The adjoint translation invariance of the model space
identifies \(\mathcal R_\alpha\) with the actual
restricted resolvent \((1-\partial)^{-1}P_\alpha\).
This is also obtained directly by integrating left
translation against \(e^{-t}\) on \(t>0\).
Equation (TB.119) gives uniformly bounded
Hilbert--Schmidt norm of \(\mathcal R_\alpha\).
Moreover, adjunction and the basis-column sum give
\[
 \begin{split}
 \|(1-P_{N,\alpha})\mathcal R_\alpha\|_2^2
 &=\|T_{r,\alpha}(P_\alpha-P_{N,\alpha})\|_2^2\\
 &\leq C_{\alpha_0}
                 \sum_{j>N}(1+\gamma_j^2)^{-1}.
 \end{split}
 \tag{TB.120}
\]
Multiplying by the uniformly Hilbert--Schmidt
\(\mathcal R_\alpha\) proves
\[
 \|(1-P_{N,\alpha})\mathcal R_\alpha^2\|_1
       \leq C_\chi
             \left(\sum_{j>N}(1+\gamma_j^2)^{-1}\right)^{1/2}.
 \tag{TB.121}
\]

We next retain the actual half-line endpoint terms.
For \(u\) in the second graph domain of left
translation, two integrations by parts give
\[
 C_k^+(1-\partial)^2u
   =C_{k-2k'+k''}^+u+(2k-k')u(0)-k\,u'(0).
 \tag{TB.122}
\]
The resolvent square maps into this domain; this can
be verified from the resolvent identity and one-dimensional
Sobolev traces, exactly as in MA-GP.96--104.
Thus (TB.122) applies to \(u=\mathcal R_\alpha^2 f\).
For \(R>b\) its endpoint terms vanish after multiplication
by \(1_{[R,\infty)}\).
Setting \(d_\alpha=k_\alpha-2k_\alpha'+k_\alpha''\)
therefore gives the actual operator identity
\[
 1_{[R,\infty)}C_{k_\alpha}^+P_\alpha
      =1_{[R,\infty)}C_{d_\alpha}^+
                                     \mathcal R_\alpha^2.
 \tag{TB.123}
\]
The norms \(\|d_\alpha\|_1\) are bounded uniformly for
\(0<\alpha\leq\alpha_0\).
Equation (TB.121) bounds the contribution of
\((1-P_{N,\alpha})\mathcal R_\alpha^2\) in trace norm.

It remains to estimate the finite model-space part
without introducing small zero gaps.
Up to phases, the inverse Laplace transform of a
normalized elementary factor is
\(\delta_0-2\alpha e^{-\overline z_{j,\alpha}t}1_{t>0}\).
The inverse transform of the initial kernel in
\(e_{j,\alpha}\) is
\(\sqrt{2\alpha}e^{-\overline z_{j,\alpha}t}1_{t>0}\).
Expanding the finite convolution of these factors and
using the volume \(t^\ell/\ell!\) of the convolution
simplex gives the pointwise estimate
\[
 |e_{j,\alpha}(t)|
       \leq\sqrt{2\alpha}\,e^{-\alpha t}p_j(\alpha t),
             \qquad t>0.
 \tag{TB.124}
\]
Every factor has the same real exponential rate
\(\alpha\); no division by a zero difference was used.
Consequently
\[
 \sum_{j=1}^N
      \|1_{[R-b,\infty)}e_{j,\alpha}\|_2^2
             \leq \mathcal H_N(\alpha(R-b))^2.
 \tag{TB.125}
\]
Since \(d_\alpha\) is supported in \([0,b]\), Young's
inequality, applied to each basis vector, yields
\[
 \|1_{[R,\infty)}C_{d_\alpha}^+P_{N,\alpha}\|_2
       \leq \|d_\alpha\|_1
                     \mathcal H_N(\alpha(R-b)).
 \tag{TB.126}
\]
Multiply by \(P_{N,\alpha}\mathcal R_\alpha^2\),
whose Hilbert--Schmidt norm is uniformly bounded.
Together with (TB.121)--(TB.123), this proves the
asserted bound for \(\omega_{\alpha,g}\).
The final inequality in (TB.116) follows by integrating
the finite polynomial times \(e^{-2v}\).

Finally, the actual Hankel identity MA-SB.4 is
\[
 H_{B_\alpha k_\alpha}
    =C_{k_\alpha}^+P_\alpha H_{B_\alpha}
                                      +H_{k_\alpha}V_{\alpha,-},
 \tag{TB.127}
\]
where both remaining blocks have norm at most one.
The second term has output support in \([0,b]\).
For \(R>b\), it follows that
\(\nu_{\alpha,g}(R)\leq\omega_{\alpha,g}(R)\);
this proves (TB.115).

The exact trace estimate MA-SB.15 applies at every
positive \(\alpha\), with right side bounded by
\(\omega_{\alpha,g}(\log(X/2))
 +\nu_{\alpha,g}(\log(X/2))\).
First choose \(N\) so that the sum in (TB.115) is
small; then (TB.117) makes its finite remainder
\(\mathcal H_N\) tend to zero.
Under \(\mathsf A\), the actual causal quotient trace
MA-GP.95 includes all zeros, and
\[
 \int_0^\infty k_\alpha(t)
                       e^{-z_{j,\alpha}t}\,dt
                =\int_0^\infty g(t)e^{i\gamma_jt}\,dt.
 \]
The series is absolutely convergent by rapid Fourier
decay and the zero count.
It is independent of \(\alpha\), proving (TB.118).
\end{proof}

\begin{remark}[The precise comparison still needed after AC-T36--T37]
These results remove the cutoff-ratio mismatch and,
for causal tests, supply a uniform shrinking-strip
estimate for the actual source PW model.
Choosing \(\alpha_\Lambda=(\log\Lambda)^{-1/2}\)
keeps the metric factors in (TB.112) subpower.
Those factors need not tend to one:
the direct small-norm metric perturbation condition
\(\alpha_\Lambda\log\max(X,Y)\to0\) is not supplied
by (TB.117).
Thus the estimates do not identify the weighted
completed PW projection with the polar projection of
the selected finite prolate theta images.
An actual tested comparison of those two projections
is still required before (TB.118) can imply the upper
direction of (TB.53).
The extension of the present uniform estimate to the
full contact-vanishing test class, and control of the
remaining identity contact for the original trace,
have not been asserted here.
\end{remark}

